%! Author = Omar Iskandarani
%! Date = 4/10/2026
%! Affiliation = Independent Researcher, Groningen, The Netherlands
%! ORCID = 0009-0006-1686-3961
%! DOI = 10.5281/zenodo.19655881
%! v0.8.1 edition: branched from v0.8.0 with dark-sector/galactic canonization integration in the main text.
%! Mint a new Zenodo DOI when publishing this edition; the v0.8.0 deposit remains \texttt{10.5281/zenodo.19682949}.
%! Build: with \texttt{-output-directory=...}, run \texttt{pdflatex} twice (or use \texttt{latexmk}) so citations and cross-refs settle.

\newcommand{\paperdoi}{10.5281/zenodo.19682949}
\newcommand{\papertitle}{Swirl-String-Theory Canon-v0.8.1}
\newcommand{\paperkeywords}{Swirl String Theory, Canon, Theoretical Physics, Formal Systems}

\documentclass[11pt]{article}
% ---------------------------
% Packages
% ---------------------------
\usepackage{amsmath,amssymb,amsfonts,bm,amsthm}
\usepackage{mathrsfs}
\usepackage{siunitx}
\usepackage[a4paper,margin=1in]{geometry}
\usepackage[absolute,overlay]{textpos}
\usepackage[T1]{fontenc}
\usepackage[utf8]{inputenc}
\usepackage{graphicx}
\usepackage{physics}
\usepackage{mathtools}
\usepackage{array}
\usepackage{import}
\usepackage{subfiles}
\usepackage{lmodern}
\usepackage{booktabs}
\usepackage{tikz}
\usetikzlibrary{arrows.meta,positioning,calc,fit,decorations.pathmorphing}
% Tables and Figures
\usepackage{float}

% ---------------------------
% Theorem environments
% ---------------------------
\theoremstyle{definition}
\newtheorem{definition}{Definition}[section]

\theoremstyle{plain}
\newtheorem{axiom}{Axiom}[section]
\newtheorem{theorem}{Theorem}[section]
\newtheorem{lemma}{Lemma}[section]

\theoremstyle{remark}
\newtheorem{remark}{Remark}[section]

% hyperref: load after float/tikz/amsthm; *before* tocloft (required by tocloft.pdf — otherwise TOC can be empty).
\usepackage{hyperref}
\hypersetup{colorlinks=true,linkcolor=blue,citecolor=blue,urlcolor=blue}

% TOC (tocloft must follow hyperref)
\usepackage{tocloft}
\setcounter{tocdepth}{2}
\renewcommand{\cftsecfont}{\footnotesize}
\renewcommand{\cftsubsecfont}{\footnotesize\itshape}
\renewcommand{\cftsecleader}{\cftdotfill{5}}
\newcommand{\compacttoc}{%
    \begingroup
    \setlength{\cftparskip}{0pt}%
    \setlength{\cftbeforesecskip}{0.5pt}%
    \setlength{\cftbeforesubsecskip}{0.5pt}%
    \renewcommand{\cftsecfont}{\scriptsize}%
    \renewcommand{\cftsubsecfont}{\scriptsize\itshape}%
    \renewcommand{\cftsecpagefont}{\scriptsize}%
    \renewcommand{\cftsubsecpagefont}{\scriptsize}%
    \setlength{\cftsecnumwidth}{2.8em}%
    \setlength{\cftsubsecindent}{1.4em}%
    \renewcommand{\cftdotsep}{2}%
    \setlength{\parskip}{0pt}%
    \linespread{0.80}\selectfont
    \tableofcontents
    \endgroup
}

% ---------------------------
% SST macro prelude
% ---------------------------
\newcommand{\swirlarrow}{\mathchoice{\mkern-2mu\scriptstyle\boldsymbol{\circlearrowleft}}{\mkern-2mu\scriptstyle\boldsymbol{\circlearrowleft}}{\mkern-2mu\scriptscriptstyle\boldsymbol{\circlearrowleft}}{\mkern-2mu\scriptscriptstyle\boldsymbol{\circlearrowleft}}}
\newcommand{\swirlarrowcw}{\mathchoice{\mkern-2mu\scriptstyle\boldsymbol{\circlearrowright}}{\mkern-2mu\scriptstyle\boldsymbol{\circlearrowright}}{\mkern-2mu\scriptscriptstyle\boldsymbol{\circlearrowright}}{\mkern-2mu\scriptscriptstyle\boldsymbol{\circlearrowright}}}
\newcommand{\vswirl}{\mathbf{v}_{\swirlarrow}}
\newcommand{\vswirlcw}{\mathbf{v}_{\swirlarrowcw}}
\newcommand{\rhoF}{{\rho_{\!f}}}
\newcommand{\rhoE}{\rho_{\!E}}
\newcommand{\rhoM}{\rho_{\!m}}
\newcommand{\rhocore}{\rho_{\mathrm{core}}}
\newcommand{\SwirlClock}{{S_{(t)}^{\swirlarrow}}}
\newcommand{\SwirlClockcw}{{S_{(t)}^{\swirlarrowcw}}}
\newcommand{\rc}{{r_c}}
\newcommand{\vnorm}{\lVert \mathbf{v}_{\!\boldsymbol{\circlearrowleft}} \rVert}



\begin{document}
    
    \thispagestyle{empty}

    % -------------------------------------------------
    % Title page / front matter
    % -------------------------------------------------
    \begin{center}
        \vspace*{0.06\textheight}

        {\LARGE\bfseries \papertitle\par}
        \vspace{0.45cm}

        {\large Canonical Reference and Research Framework\par}
        \vspace{1.0cm}

        {\Large Omar Iskandarani\textsuperscript{*}\par}
        \vspace{0.35cm}

        {\normalsize Independent Researcher\par}
        {\normalsize Groningen, The Netherlands\par}
        \vspace{0.6cm}

        {\normalsize Version v0.8.1\par}
        {\normalsize \today\par}
        \vspace{0.9cm}

        \rule{0.82\textwidth}{0.5pt}
        \vspace{0.7cm}
    \end{center}

    \begin{abstract}
        \noindent
        This document presents Swirl-String Theory (SST) as a structured canonical synthesis rather than as a claim that every phenomenological sector has already been derived to theorem level. The text consolidates the primitive constant chain linking the effective fluid density, characteristic swirl speed, Compton anchoring, and core scale; records delay-induced mode selection as the preferred non-operator route to spectral discreteness; and organizes the atomic bridge, spectroscopic filter, relational-time sector, and topology-based mass program within one explicitly stratified framework. Throughout, orthodox inputs, SST-internal derivations, bridge-level constructions, and conjectural extensions are kept distinct so that the canon remains internally coherent, externally comparable, and open to falsification.
    \end{abstract}

    \vspace{0.2cm}
    \noindent\textbf{Keywords:} \paperkeywords

    \vspace{0.5cm}
    \noindent\rule{\textwidth}{0.4pt}
    \vspace{0.35cm}
    
    \noindent\begin{quote}
        \small
        \textbf{Canon reading guide.}
        
        \vspace{0.3em}
        \begin{tabular}{@{}>{\raggedleft\arraybackslash}p{0.21\textwidth}@{\hspace{0.32em}}>{\raggedright\arraybackslash}p{0.75\textwidth}@{}}
            \textbf{[ORTHODOX]} & established physics or mathematics. \\[0.25em]
            \textbf{[DERIVED]} & consequences obtained from definitions and assumptions of this SST layer. \\[0.25em]
            \textbf{[SPECULATIVE]} & conjectural extensions or identifications not yet closed as theorems. \\
        \end{tabular}
        
        \vspace{0.4em}
        Optional \emph{role} labels (\emph{e.g.}, \emph{Definition}, \emph{Benchmark}, \emph{Constraint}) qualify how to read a statement but do not replace the primary tag.
        Orthodox constraints---including quantum predictions, atomic spectroscopy, and relativistic consistency---are not overridden: bridge constructions, benchmark sections, and research-track material remain subordinate to the consistency stratification introduced later in the canon.
    \end{quote}

    \medskip

    \begin{textblock*}{1\textwidth}(0.15\textwidth,1.1\textheight)\raggedright\footnotesize\renewcommand{\arraystretch}{1.0} \noindent\rule{\textwidth}{0.4pt}\\[0.5em]
        \textsuperscript{\textbf{*}} Email: \texttt{info@omariskandarani.com}\\
        ORCID: \texttt{\href{https://orcid.org/0009-0006-1686-3961}{0009-0006-1686-3961}}\\
        DOI: \href{https://doi.org/\paperdoi}{\paperdoi}
    \end{textblock*}
    
    \newpage
    % Original: \tableofcontents
    % Use compact version to save space
    \compacttoc
    \newpage


% ===========================
% CANON STRUCTURE (v0.8.1)
% ===========================

% ---------- PRELUDE ----------
    \section{Philosophical Prelude}
        \label{sec:philosophy}

    \subsection{Why a philosophical prelude is necessary}
        The SST canon is not only a collection of equations but also a statement about what counts as fundamental. Standard quantum theory and general relativity are highly successful formalisms, yet they begin from different ontological primitives: operator-valued states in one case and spacetime geometry in the other. SST instead adopts a single material substrate as the underlying level of description and treats particles, radiation, and clock structure as organized states of that substrate.

    \subsection{The ontological shift}
        The canon therefore makes the following conceptual replacement:
        \begin{align}
            \text{point particles} &\longrightarrow \text{stable swirl-string configurations}, \\
            \text{vacuum as background} &\longrightarrow \text{vacuum as structured medium}, \\
            \text{quantization as postulate} &\longrightarrow \text{discreteness from dynamics and topology}.
        \end{align}
        This is not a denial of effective particle language. Rather, it is a claim that particle language is secondary and emerges from a more primitive hydrodynamic and topological layer.

    \subsection{Relation to orthodox physics}
        SST is not introduced as a replacement for the empirical content of orthodox physics. Its canonical role is instead threefold:
        \begin{enumerate}
            \item to reproduce known low-energy structures in the appropriate limits,
            \item to identify which apparently fundamental constants or rules may be structurally redundant,
            \item to supply a dynamical mechanism for phenomena that are otherwise postulated axiomatically.
        \end{enumerate}
        In this sense the canon follows a \emph{Rosetta principle}: whenever a standard description already works, SST must map to it before claiming any deeper reinterpretation.

    \subsection{Why delay and topology matter}
        Two themes recur across the development of the canon.
        First, finite circulation delay on closed carriers provides a natural source of temporal nonlocality. Second, topological invariants constrain admissible configurations and stabilize structures that would otherwise collapse or disperse. Together they motivate the canonical SST slogan
        \begin{center}
            \textit{delay selects modes, topology protects them.}
        \end{center}
        This is the philosophical reason why the canon treats quantization neither as a primitive operator rule nor as a purely statistical statement.

    \subsection{Scope of the present canon}
        The present version of the canon should be read as a structured reference document rather than a claim that every sector has already been derived to theorem level. Some components are mature and reusable, such as the primitive constant chain, circulation quantization, and the Swirl-Clock. Other components are deliberately retained as research-track layers, such as full gauge emergence, topology-to-charge assignments, and complete many-body atomic closure. The point of the canon is therefore not maximal rhetorical closure, but explicit logical organization.

    \subsection{Interpretive summary}
        The philosophical stance of SST can be summarized in one sentence:
        \begin{center}
            \textit{Physics is interpreted as the dynamics of a continuous medium whose stable, delayed, and topologically organized motions appear effectively as particles, clocks, fields, and spectra.}
        \end{center}

% ---------- FOUNDATIONS ----------
\section{Formal Foundations}
    \label{sec:foundations}
    
    \subsection{Primitive Structure of the Theory}
        
        We define the Swirl-String Theory (SST) as a continuum-based framework built on a minimal set of primitive quantities.
        
        \textbf{Primitive set:}
        \begin{align}
            \mathcal{P} = \{ \rhoF, \vswirl, \omega_c \}
        \end{align}
        
        where:
        \begin{itemize}
            \item $\rhoF$ is the effective background fluid density,
            \item $\vswirl$ is the characteristic swirl velocity,
            \item $\omega_c$ is the Compton angular frequency.
        \end{itemize}
        
        \textbf{[ORTHODOX]}: $\omega_c = \frac{m_e c^2}{\hbar}$ is a standard physical constant.
        
        \textbf{[SPECULATIVE]}: The identification of $\omega_c$ as a primitive dynamical frequency of the vortex core is an SST postulate.
    
    \subsection{Derived Length Scale}
        
        The characteristic core radius is defined as:
        
        \begin{align}
            \rc = \frac{\vnorm}{\omega_c}
            \label{eq:rc_definition}
        \end{align}
        
        \textbf{[SPECULATIVE] Definition}: We treat Eq.~\eqref{eq:rc_definition} as a canonical closure relation linking the characteristic core scale to the primitive pair $(\vnorm,\omega_c)$.

        \textbf{[DERIVED] Consequence}: This removes $r_c$ as an independent parameter and eliminates one source of circular parameterization.
    
    \subsection{Circulation Quantization}
        
        The circulation is defined as:
        
        \begin{align}
            \Gamma = \oint \mathbf{v} \cdot d\ell
        \end{align}
        
        We introduce the fundamental circulation quantum:
        
        \begin{align}
            \Gamma_0 = 2\pi \rc \, \vnorm
            \label{eq:gamma0}
        \end{align}
        
        and impose:
        
        \begin{align}
            \Gamma = n \Gamma_0, \quad n \in \mathbb{Z}
        \end{align}
        
        \textbf{[ORTHODOX]}: Circulation quantization appears in superfluid systems.
        
        \textbf{[DERIVED]}: In SST, $\Gamma_0$ is no longer an input but follows from Eq.~\ref{eq:rc_definition}.
    
    \subsection{Energy and Mass Relation}
        
        Mass is interpreted as stored rotational energy:
        
        \begin{align}
            E = m c^2
        \end{align}
        
        \textbf{[ORTHODOX]}: Standard relativistic energy relation.
        
        \textbf{[SPECULATIVE]}: SST interprets this energy as localized vortex kinetic energy.
    
    \subsection{Swirl Velocity Constraint}
        
        We define the dimensionless ratio:
        
        \begin{align}
            \eta_0 = \frac{\vnorm}{c}
        \end{align}
        
        Empirically:
        
        \begin{align}
            \eta_0 \approx \frac{\alpha}{2}
        \end{align}
        
        \textbf{[DERIVED] Empirical closure}: Within the present canonical calibration, this connects the characteristic swirl velocity to the fine-structure constant.

        \textbf{[CRITICAL NOTE]} This relation should be read as a calibrated SST closure unless and until it is re-derived from a deeper dynamical sector.
    
    \subsection{Swirl-Clock Definition}
        
        We introduce the Swirl-Clock as a local scalar field encoding relativistic time dilation:
        
        \begin{align}
            \SwirlClock = \sqrt{1 - \frac{\vnorm^2}{c^2}}
            \label{eq:swirl_clock}
        \end{align}
        
        \textbf{[ORTHODOX] Functional form}: This matches the standard Lorentz time-dilation factor.

        \textbf{[SPECULATIVE] Identification}: In SST, the relevant velocity is interpreted as a local medium-kinematic state rather than only a standard worldline speed.

        \textbf{[DERIVED] Consequence}: Once this identification is adopted, the Swirl-Clock provides a local scalar measure of kinematic time modulation within the medium description.
    
    \subsection{Density Hierarchy}
        
        We distinguish the following densities:
        
        \begin{align}
            \rhoF &\quad \text{(effective background fluid density)}, \\
            \rhoE &= \tfrac12 \rhoF \, \vnorm^2 \quad \text{(swirl energy density)}, \\
            \rhoM &= \rhoE / c^2 \quad \text{(mass-equivalent density)}, \\
            \rhocore &\quad \text{(saturated vortex-core density)}.
        \end{align}

        \textbf{[ORTHODOX] Definition}: $\rhoM=\rhoE/c^2$ is the mass-equivalent density associated with the swirl kinetic accounting above; in orthodox limits, total energy $E$ in a region still yields mass density $E/c^2$.

        \textbf{[SPECULATIVE] Interpretation}: The distinction between background, energetic, and saturated-core sectors is elevated in SST to a structural ontology of the medium.
    
    \subsection{Core Density Closure}
        
        The core density follows from a force constraint:
        
        \begin{align}
            \rhocore \sim \frac{m_e c^2}{2\pi \vnorm^2 \rc^3}
            \label{eq:core_density}
        \end{align}
        
        \textbf{[SPECULATIVE] Closure ansatz}: Equation~\eqref{eq:core_density} is retained as the canonical core-density closure used by the present SST calibration program.

        \textbf{[DERIVED] Consequence}: Within that closure, $\rhocore$ is no longer an independent parameter.

    \subsection{Chronos--Kelvin Invariant and Clock--Radius Transport}
        \label{subsec:chronos_kelvin}

        For a thin material loop of radius $R(t)$ carried by an incompressible, inviscid flow with no reconnection events, Kelvin's theorem implies conservation of circulation. In the near-solid-body core approximation this gives

        \begin{align}
            \frac{D}{Dt}\!\bigl(R^2 \omega\bigr) = 0,
            \label{eq:chronos_kelvin_basic}
        \end{align}

        where $\omega$ is the characteristic local vorticity on the loop. Using the core relation $v_t = \omega \rc$ together with Eq.~\eqref{eq:swirl_clock}, one may rewrite

        \begin{align}
            \omega = \frac{c}{\rc}\sqrt{1-S_{(t)}^2},
        \end{align}

        so that Eq.~\eqref{eq:chronos_kelvin_basic} takes the equivalent form

        \begin{align}
            \frac{D}{Dt}\left(\frac{c}{\rc} R^2 \sqrt{1-S_{(t)}^2}\right) = 0.
            \label{eq:chronos_kelvin_clock}
        \end{align}

        \textbf{[ORTHODOX]}: Eq.~\eqref{eq:chronos_kelvin_basic} is the thin-loop form of Kelvin circulation conservation in inviscid barotropic flow \cite{Batchelor1967,Saffman1992}.

        \textbf{[DERIVED]}: Eq.~\eqref{eq:chronos_kelvin_clock} is the SST rewriting of the same invariant in terms of the Swirl-Clock.

        Differentiating Eq.~\eqref{eq:chronos_kelvin_clock} gives a useful transport identity:

        \begin{align}
            \frac{dS_{(t)}}{dt}
            =
            \frac{2\bigl(1-S_{(t)}^2\bigr)}{S_{(t)}}
            \frac{1}{R}\frac{dR}{dt}.
            \label{eq:clock_radius_transport}
        \end{align}

        \textbf{[DERIVED]}: contraction of a material loop drives stronger clock suppression, while expansion relaxes it.

    \subsection{Reparametrized Zero-Parameter Corollary}
        \label{subsec:zero_parameter_corollary}

        Older SST canon layers often used the triplet $(\Gamma_0,\rhoF,\rc)$ as the primitive parameterization. In the present v0.8.1 architecture, Eq.~\eqref{eq:rc_definition} and Eq.~\eqref{eq:gamma0} show that this is a reparameterization of the present primitive set $(\rhoF,\vswirl,\omega_c)$ rather than a distinct axiom.

        Explicitly,

        \begin{align}
            \Gamma_0 = 2\pi \rc \, \vnorm
            = 2\pi \frac{\vnorm^2}{\omega_c}.
            \label{eq:gamma0_reparam}
        \end{align}

        Hence any dimensional SST quantity written in terms of $(\Gamma_0,\rhoF,\rc)$ can be rewritten in terms of $(\rhoF,\vswirl,\omega_c)$ with no new input.

        A coarse-grained density law retained from earlier canon layers is

        \begin{align}
            \rhoF = K\,\Omega,
            \qquad
            K = \frac{\rhocore \, \rc}{\vnorm},
            \label{eq:coarse_grain_density}
        \end{align}

        where $\Omega$ is a coarse-grained local rotation rate and $K$ has dimensions $\mathrm{kg\,m^{-3}\,s}$. This relation records the canonical bridge between a sparse effective bulk density and a high-density localized core.

        \textbf{[DERIVED COROLLARY]}: the older ``zero-parameter'' language is retained only as a statement of reparameterization discipline: once one primitive triplet is fixed, no additional dimensional constants should be introduced without explicit justification.

    \subsection{Canonical Master Equation}
    \label{sec:master-equation}

        The preceding subsections define the four ingredients that later SST papers often treat separately:
        a local mass-equivalent density, a local Swirl-Clock, a geometric gate, and a topological kernel.
        We therefore introduce a single canonical equation that organizes these sectors in one place.

        First define the local swirl energy density and its mass-equivalent form:

        \begin{align}
            \rhoE(x) &= \frac{1}{2}\rhoF(x)\,\norm{\vswirl(x)}^2 \\
            \rhoM(x) &= \frac{\rhoE(x)}{c^2}
        \end{align}

        Using Eq.~\ref{eq:swirl_clock}, the local Swirl-Clock factor is

        \begin{align}
            \SwirlClock(x) = \sqrt{1 - \frac{\norm{\vswirl(x)}^2}{c^2}}
        \end{align}

        We then define a binary geometric exposure gate
        \begin{align}
            G(K) \in \{0,1\}
        \end{align}
        together with the canonical geometric factor
        \begin{align}
            \Pi_K = \left( \frac{\lambda_c}{\pi \rc} \right)^{G(K)} = \left( \frac{4}{\alpha} \right)^{G(K)} .
        \end{align}

        The canonical topological kernel is taken to be
        \begin{align}
            \Xi_K = \left[ \alpha_C\,C(K) + \beta_L\,L(K) + \gamma_H\,\mathcal{H}(K) \right] \varphi^{-2k(K)} .
        \end{align}

        The canonical SST master equation is therefore
        \begin{align}
            M_K(x)
            &= \rhoM(x)\,\rc^3\,\Pi_K\,\Xi_K\,\SwirlClock(x)^{-2} \\
            &= \left( \frac{\rhoF(x)\,\norm{\vswirl(x)}^2}{2c^2} \right)
               \rc^3
               \left( \frac{\lambda_c}{\pi \rc} \right)^{G(K)}
               \left[ \alpha_C\,C(K) + \beta_L\,L(K) + \gamma_H\,\mathcal{H}(K) \right]
               \varphi^{-2k(K)}
               \SwirlClock(x)^{-2}
            \label{eq:sst_master_equation}
        \end{align}

        \textbf{[SPECULATIVE] Organizing principle}: This equation factorizes SST into four sectors:
        \begin{align}
            \text{medium scale} \times \text{geometric gate} \times \text{topological kernel} \times \text{clock impedance} .
        \end{align}

        \textbf{[DERIVED] Dimensional check}:
        \begin{align}
            \left[ \frac{\rhoF\,\norm{\vswirl}^2}{c^2} \rc^3 \right] = \mathrm{kg}
        \end{align}
        and all remaining factors are dimensionless, so Eq.~\ref{eq:sst_master_equation} has the correct mass dimension.

        \textbf{[SPECULATIVE] Interpretation}: The same local kinematic state that slows the Swirl-Clock also enters the effective mass scale; mass and local proper-time modulation are therefore treated as two aspects of one medium configuration.

        \textbf{[CRITICAL NOTE]} Equation~\eqref{eq:sst_master_equation} is presently a canonical synthesis formula, not yet a theorem derived from first-principles filament dynamics alone.

        In the stationary weak-swirl limit $\SwirlClock \to 1$, it reduces to the static mass-functional form; in the local-kinematic limit at fixed $K$, it reduces to a pure clock-impedance correction. The hierarchy is therefore carried predominantly by the geometric and topological factors, while the Swirl-Clock supplies a coherent local correction rather than the primary source of mass separation.

    \subsection{Geometric Representation}
        
        A swirl-string is represented as a closed curve:
        
        \begin{align}
            X(t,\sigma) = (x(t,\sigma), y(t,\sigma), z(t,\sigma))
        \end{align}
        
        with $\sigma \in [0,2\pi)$.
    
    \subsection{Summary of Dependency Structure}
        
        The full dependency chain is:
        
        \begin{align}
        (\rhoF, \vswirl, \omega_c)
            \rightarrow \rc
            \rightarrow \Gamma_0
            \rightarrow (\rhoE, \rhoM, \rhocore)
            \rightarrow \SwirlClock
            \rightarrow M_K.
        \end{align}
        
        \textbf{[DERIVED] Consistency statement}:
        No quantity is defined circularly. The canonical closure is summarized by Eq.~\ref{eq:sst_master_equation}.
    
    \subsection{Consistency Statement}
        
        \begin{itemize}
            \item All derived quantities depend only on $\mathcal{P}$
            \item No external parameters are introduced
            \item The theory is closed under its primitive set
        \end{itemize}
    
    \subsection{Interpretation}
        
        The Formal Foundations establish SST as:
        
        \begin{itemize}
            \item a continuum theory with minimal primitives
            \item a geometrically grounded framework
            \item a non-circular parameter system
        \end{itemize}

    \subsection{Maximal Swirl Tension}
    \label{sec:Fmax}

        \subsubsection{Definition and canonical form}

        Within the primitive set $\mathcal{P} = \{\rhoF, \vswirl, \omega_c\}$, a characteristic force scale emerges naturally from the Compton energy stored over the core cross-section. We define the \emph{maximal swirl tension} as

        \begin{align}
            F_{\rm swirl}^{\max}
            \;=\; \frac{\vnorm\,\hbar}{2\rc^2}
            \;=\; \frac{\hbar\omega_c}{2\rc}.
            \label{eq:Fmax_def}
        \end{align}

        \textbf{[DERIVED] Canonical form}: substituting $\vnorm = \omega_c \rc$ into Eq.~\eqref{eq:Fmax_def} gives $F_{\rm swirl}^{\max} = \hbar\omega_c/2\rc$, the Compton energy scale divided by twice the core radius. Numerically, $F_{\rm swirl}^{\max} = 29.053507\;\text{N}$.

        \textbf{[CRITICAL NOTE]} $F_{\rm swirl}^{\max}$ is a mechanical reformulation of the Compton energy scale over the core cross-section. It is \emph{not} independent of $\hbar$ and $m_e$; its role in the canon is interpretive and organisational rather than that of an independent primitive.

        \subsubsection{Internal consistency identities}

        The following expressions are algebraically equivalent within the canonical constant chain and serve as coherence verifications rather than independent derivations:

        \begin{align}
            F_{\rm swirl}^{\max}
            &= \frac{e^2}{16\pi\varepsilon_0 \rc^2},
            \label{eq:Fmax_em}\\
            &= \frac{h\alpha c}{8\pi \rc^2},
            \label{eq:Fmax_hac}\\
            &= \pi \rc^2\,\rho_{\rm calc}\,\vnorm^2,
            \label{eq:Fmax_rho}
        \end{align}

        where $\rho_{\rm calc} = 4F_{\rm swirl}^{\max}/(\pi\alpha^2 c^2 \rc^2)$ is the derived bulk density consistent with the canonical core closure.

        \textbf{[DERIVED] Consistency check}: Eqs.~\eqref{eq:Fmax_em}--\eqref{eq:Fmax_rho} are equivalent given $\vnorm = \alpha c/2$ and the standard identity $e^2/(4\pi\varepsilon_0) = \alpha\hbar c$. Their mutual agreement constitutes a non-trivial numerical self-consistency test of the primitive set.

        \subsubsection{Strongest non-circular identity: the Rydberg constant}

        \textbf{[DERIVED]} The Rydberg constant is recoverable directly from the three SST primitives without $e$, $\varepsilon_0$, or $m_e$ as explicit input:

        \begin{align}
            R_\infty \;=\; \frac{\vnorm^3}{\pi\,\rc\,c^3}.
            \label{eq:Rydberg_SST}
        \end{align}

        This is the canonical benchmark for the internal non-redundancy of the primitive set. Combining Eq.~\eqref{eq:Rydberg_SST} with Eq.~\eqref{eq:Fmax_def} yields the compact relation

        \begin{align}
            \boxed{
            F_{\rm swirl}^{\max}
            \;=\; \frac{32\pi^2\hbar R_\infty^2 c}{\alpha^5},
            }
            \label{eq:Fmax_Rydberg}
        \end{align}

        which expresses $F_{\rm swirl}^{\max}$ in terms of spectroscopic observables alone, with no reference to $m_e$, $\rc$, or $e$ as separate inputs. A further identity connecting atomic and force scales is

        \begin{align}
            R_\infty \;=\; F_{\rm swirl}^{\max} \cdot \frac{\alpha^3}{4\pi m_e c^2},
            \label{eq:Rydberg_Fmax}
        \end{align}

        which states that the Rydberg constant is the maximal swirl tension scaled by the Compton energy and a purely geometric factor $\alpha^3/4\pi$.

        \subsubsection{Threefold physical interpretation}

        The quantity $F_{\rm swirl}^{\max}$ admits three distinct physical interpretations, each arising from a different sector of physics. Their numerical coincidence within the canonical constant chain is a non-trivial coherence result.

        \paragraph{(i) Gravitational analogue: the Planck force ceiling.}

        In general relativity, the combination $c^4/4G$ sets the maximum force that any physical process can transmit without forming a causal horizon~\cite{Gibbons2002}. We denote this scale the maximal gravitational tension:

        \begin{align}
            F_{\rm gr}^{\max} \;=\; \frac{c^4}{4G}.
            \label{eq:Fgr}
        \end{align}

        \textbf{[ORTHODOX]} $F_{\rm gr}^{\max}$ is the Planck force, an invariant upper bound in classical general relativity \cite{Gibbons2002}.

        \textbf{[DERIVED]} The ratio of the two tension scales reduces exactly to

        \begin{align}
            \boxed{
            \frac{F_{\rm swirl}^{\max}}{F_{\rm gr}^{\max}}
            \;=\; \frac{4\alpha_g}{\alpha},
            }
            \label{eq:Fmax_ratio}
        \end{align}

        where $\alpha_g = Gm_e^2/\hbar c$ is the gravitational fine-structure constant of the electron and $\alpha$ is the electromagnetic fine-structure constant. Numerically, $4\alpha_g/\alpha \approx 9.6\times10^{-43}$, in agreement with the directly computed ratio $29.053507\;\text{N}\,/\,3.026\times10^{43}\;\text{N}$.

        \textbf{[SPECULATIVE] Interpretation}: the hierarchy between the electromagnetic and gravitational force scales in SST is not an unexplained numerical coincidence but the ratio of the two fundamental coupling constants of the electron. The structured medium is $\sim\!10^{42}$ times stiffer than the space-time geometry at the core scale, and this stiffness ratio equals $\alpha/(4\alpha_g)$ exactly.

        \paragraph{(ii) Electrostatic interpretation: the Coulomb ceiling at core scale.}

        When two protons are compressed to a centre-to-centre separation equal to the vortex core radius $\rc$, the Coulomb repulsion is

        \begin{align}
            F_{pp}(\rc)
            \;=\; \frac{e^2}{4\pi\varepsilon_0 \rc^2}
            \;=\; 4\,F_{\rm swirl}^{\max}.
            \label{eq:Coulomb_core}
        \end{align}

        \textbf{[DERIVED]} The maximal Coulomb repulsion at core scale is exactly four times $F_{\rm swirl}^{\max}$. The factor of four matches the denominator structure of Eq.~\eqref{eq:Fgr} and is not numerically coincidental within the canonical chain.

        \textbf{[SPECULATIVE] Interpretation}: $F_{\rm swirl}^{\max}$ is the quarter-Coulomb force at core scale. The medium sustains one quarter of the maximum proton--proton Coulomb repulsion before the topological structure of the vortex core fails. This provides a physical ceiling on the compression of charged vortex configurations.

        \paragraph{(iii) Mechanical interpretation: the spring constant of photon--electron interaction.}

        The electron's internal vortex is modelled as a linear spring with stiffness

        \begin{align}
            K_e \;=\; \frac{F_{\rm swirl}^{\max}}{n\,\rc},
            \qquad n = 2,
            \label{eq:spring_K}
        \end{align}

        where $n$ counts the number of half-wavelength segments on the core and is fixed to $n=2$ by the photon--electron swirl matching condition. The natural oscillation frequency of the electron core under this spring is

        \begin{align}
            \omega_{\rm spring}
            \;=\; \sqrt{\frac{K_e}{m_e}}
            \;=\; \sqrt{\frac{F_{\rm swirl}^{\max}}{2\,\rc\,m_e}}
            \;=\; \frac{\omega_c}{\alpha}
            \;=\; \frac{m_e c^2}{\alpha\hbar},
            \label{eq:omega_spring}
        \end{align}

        a frequency intermediate between the atomic Rydberg scale $(\sim\alpha^2\omega_c)$ and the bare Compton scale $(\omega_c)$.

        \textbf{[DERIVED]} The energy stored in this spring at maximum compression $\delta = \rc$ is

        \begin{align}
            E_{\rm spring}
            \;=\; \tfrac{1}{2}\,K_e\,\rc^2
            \;=\; \frac{F_{\rm swirl}^{\max}\,\rc}{4}
            \;=\; \frac{\hbar\omega_c}{8}
            \;=\; \frac{m_e c^2}{8}.
            \label{eq:Espring}
        \end{align}

        One eighth of the electron rest energy is stored as internal vortex-spring potential at maximum compression.

        \textbf{[SPECULATIVE] Physical picture}: during photon absorption, the electron vortex core is compressed for a duration of order one Planck time $t_p$ into a transient mass-like configuration. The restoring force scale of this compression is $F_{\rm swirl}^{\max}$, and the energy scale of the transient is $m_e c^2/8$. The spring releases at the Compton period, generating the recoil that initiates the Coulomb interaction with the proton.

        \textbf{[CRITICAL NOTE]} Equation~\eqref{eq:omega_spring} assumes $n=2$ from empirical matching. A first-principles derivation of the photon wrap number from the delay dynamics of Section~\ref{sec:delay} is an open derivational step.

        \subsubsection{Unified summary}

        The three interpretations are collected in Table~\ref{tab:Fmax_triple}. Together they identify $F_{\rm swirl}^{\max}$ as the \emph{structural tension threshold of the medium}: the scale at which a physical system --- gravitational, electrostatic, or mechanical --- exhausts its capacity to store energy in the current configuration and must either fail topologically, form a horizon, or release energy as radiation.

        \begin{table}[h]
        \centering
        \caption{Threefold physical interpretation of $F_{\rm swirl}^{\max}$ within the SST canonical framework.}
        \label{tab:Fmax_triple}
        \begin{tabular}{p{3.2cm}p{5.0cm}p{4.8cm}l}
        \hline
        \textbf{Context} & \textbf{Expression} & \textbf{Physical meaning} & \textbf{Status} \\
        \hline
        Gravitational ceiling &
          $F_{\rm gr}^{\max}/F_{\rm swirl}^{\max} = \alpha/(4\alpha_g)$ &
          Medium stiffer than spacetime by $\alpha/4\alpha_g$ &
          {[Derived]} \\[6pt]
        Coulomb ceiling &
          $F_{pp}(\rc) = 4\,F_{\rm swirl}^{\max}$ &
          Quarter of max.\ Coulomb force at core scale &
          {[Derived]} \\[6pt]
        Vortex spring &
          $K_e = F_{\rm swirl}^{\max}/2\rc$ &
          Spring constant of photon--electron interaction &
          {[Speculative]} \\
        \hline
        \end{tabular}
        \end{table}

        \textbf{[SPECULATIVE] Canonical interpretation}: $F_{\rm swirl}^{\max}$ is the tension at which the structured \ae ther medium reaches its topological integrity limit. It plays the same organisational role for the electromagnetic and mechanical sectors that $F_{\rm gr}^{\max} = c^4/4G$ plays for the gravitational sector. Their ratio is precisely the ratio of the two fundamental coupling constants of the electron.

% ---------- AXIOMS ----------
\section{Axiomatic Framework}
    \label{sec:axioms}
    
    \subsection{Preliminaries}
        
        We formalize Swirl-String Theory (SST) as an axiomatic continuum theory defined over a three-dimensional spatial domain $\mathcal{M} \subset \mathbb{R}^3$ equipped with a preferred time parameter $t \in \mathbb{R}$.
        
        The theory is constructed from a minimal primitive set $\mathcal{P}$ (Section~\ref{sec:foundations}) and a set of axioms $\mathcal{A}$ governing admissible configurations and dynamics.
    
    \subsection{Axiom 1: Continuum Substrate}
        
        \textbf{Statement.}
        There exists a continuous, incompressible, inviscid medium filling $\mathcal{M}$, characterized by a constant background density $\rhoF$.
        
        \begin{align}
            \nabla \cdot \mathbf{v} = 0
        \end{align}
        
        \textbf{[ORTHODOX]}: This corresponds to the standard incompressible Euler framework.
        
        \textbf{Interpretation.}
        All physical structures arise as configurations within this medium; no discrete particles are assumed at the fundamental level.
    
    \subsection{Axiom 2: Existence of Stable Vortex Filaments}
        
        \textbf{Statement.}
        The medium admits stable, localized, topologically nontrivial vortex filament solutions.
        
        \textbf{[ORTHODOX]} Physics input: vortex filaments are known idealized structures in fluid dynamics.

        \textbf{[SPECULATIVE]} Extension: such filaments are assumed here to be dynamically stable and persist as fundamental carriers.
        
        \textbf{Mathematical form.}
        A filament is represented as a closed embedding:
        
        \begin{align}
            X: S^1 \to \mathcal{M}, \quad \sigma \mapsto X(t,\sigma)
        \end{align}
    
    \subsection{Axiom 3: Circulation Quantization}
        
        \textbf{Statement.}
        Circulation around any stable filament is quantized:
        
        \begin{align}
            \Gamma = \oint \mathbf{v} \cdot d\ell = n \Gamma_0, \quad n \in \mathbb{Z}
        \end{align}
        
        \textbf{[ORTHODOX]}: Quantized circulation appears in superfluid systems.
        
        \textbf{[SPECULATIVE]}: Elevated here to a universal invariant.
    
    \subsection{Axiom 4: Compton Anchoring}
        
        \textbf{Statement.}
        The characteristic angular frequency of the vortex core satisfies:
        
        \begin{align}
            \frac{\vnorm}{\rc} = \omega_c
        \end{align}
        
        \textbf{[ORTHODOX]}: $\omega_c$ is the Compton frequency.
        
        \textbf{[SPECULATIVE]}: Identified as the intrinsic rotational frequency of the vortex core.
        
        \textbf{Consequence.}
        The core radius is not independent but derived (Eq.~\ref{eq:rc_definition}).
    
    \subsection{Axiom 5: Delay as a Constitutive Principle}
        
        \textbf{Statement.}
        Finite propagation time along closed vortex structures is a fundamental dynamical ingredient.
        
        \begin{align}
            \tau = \frac{L}{\vnorm}
        \end{align}
        
        \textbf{[DERIVED]} Kinematic consequence: the delay follows from finite signal speed along a closed carrier.

        \textbf{Axiomatic role.} The delay is not treated as negligible, but as structurally relevant to the dynamics.
        
        \textbf{Interpretation.}
        Temporal nonlocality is intrinsic to the dynamics.
    
    \subsection{Axiom 6: Energy Localization}
        
        \textbf{Statement.}
        Energy is localized in vortex structures and determines mass via:
        
        \begin{align}
            E = m c^2
        \end{align}
        
        \textbf{[ORTHODOX]}: Relativistic energy relation.
        
        \textbf{[SPECULATIVE]}: Interpreted as rotational energy of the vortex configuration.
    
    \subsection{Axiom 7: Swirl-Clock Relational Time}
        
        \textbf{Statement.}
        Local time is encoded in a scalar field derived from fluid kinematics:
        
        \begin{align}
            \SwirlClock = \sqrt{1 - \frac{\vnorm^2}{c^2}}
        \end{align}
        
        \textbf{[ORTHODOX]} Functional form: this matches Lorentz time dilation.

        \textbf{[SPECULATIVE]} Identification: the relevant velocity is taken to be the local medium kinematic state.

        \textbf{[DERIVED]} Consequence: within SST this yields a relational clock field tied to local motion.
        
        \textbf{Interpretation.}
        Time is not fundamental but relational to the state of motion in the medium.
    
    \subsection{Axiom 8: Density Saturation and Core Structure}
        
        \textbf{Statement.}
        The vortex core exhibits a saturated density $\rhocore$ determined by mechanical constraints of the medium.
        
        \textbf{[SPECULATIVE]} Closure hypothesis: the saturated density is fixed by the closure relation of Eq.~\eqref{eq:core_density}.

        \textbf{[DERIVED]} Consequence: once that closure is adopted, matter corresponds to regions where energy density reaches a structural limit.
    
    \subsection{Logical Structure of the Axioms}
        
        The axioms form a dependency hierarchy:
        
        \begin{align}
            \text{Continuum}
            \rightarrow \text{Vortex Structures}
            \rightarrow \text{Quantization}
            \rightarrow \text{Delay Dynamics}
            \rightarrow \text{Mass/Energy Emergence}
        \end{align}
    
    \subsection{Minimality Principle}
        
        \textbf{Statement.}
        The axioms are minimal in the sense that removing any one of them destroys at least one of the following:
        
        \begin{itemize}
            \item existence of stable structures
            \item emergence of discrete spectra
            \item compatibility with known physics
        \end{itemize}
    
    \subsection{Philosophical Interpretation}
        
        The axiomatic structure implies:
        
        \begin{itemize}
            \item Particles are not fundamental objects but stable dynamical configurations
            \item Discreteness is emergent, not postulated
            \item Time is relational, not absolute
        \end{itemize}
    
    \subsection{Consistency Condition}
        
        A valid SST realization must satisfy:
        
        \begin{itemize}
            \item Compatibility with orthodox physics in appropriate limits
            \item Internal logical closure
            \item Explicit traceability from axioms to predictions
        \end{itemize}
    
    \subsection{Canonical Statement}
        
        \begin{center}
            \textit{
                Swirl-String Theory is defined as the minimal continuum theory in which stable vortex structures, governed by quantized circulation and delay dynamics, give rise to discrete physical phenomena consistent with known physics.
            }
        \end{center}

% ---------- CONSISTENCY LAYER ----------
%    \section{Consistency and Classification Framework}
%        \label{sec:consistency}
        
\section{Consistency and Classification Framework}
    \label{sec:consistency}
    
    \subsection{Classification Scheme}
        
        All statements are labeled as:
        
        \begin{itemize}
            \item \textbf{[ORTHODOX]} — established physics
            \item \textbf{[DERIVED]} — logically deduced within SST
            \item \textbf{[SPECULATIVE]} — novel or unverified
        \end{itemize}

    \subsection{Extended Labeling Convention}

        We distinguish between \emph{epistemic status} and \emph{editorial role}.

        \paragraph{Epistemic status.}
        The only primary status labels used in the present Canon are:
        \begin{itemize}
            \item \textbf{[ORTHODOX]} --- established in standard physics or mathematics
            \item \textbf{[DERIVED]} --- obtained internally from the present SST assumptions, definitions, or closures
            \item \textbf{[SPECULATIVE]} --- conjectural SST extension, interpretation, or unverified identification
        \end{itemize}

        \paragraph{Editorial role.}
        Additional descriptors such as \emph{Definition}, \emph{Constraint}, \emph{Interpretation}, \emph{Benchmark}, \emph{Roadmap}, and \emph{Dimensional check} are not independent epistemic classes. They indicate only the role played by the statement.

        \paragraph{Usage rule.}
        When both are needed, the Canon uses the format
        \begin{align}
            \texttt{[STATUS] Role: statement.}
        \end{align}
        For example, one may write \textbf{[ORTHODOX] Constraint}, \textbf{[DERIVED] Consequence}, or \textbf{[SPECULATIVE] Interpretation}.
    
    \subsection{Orthodox Constraints}
        
        The following must remain intact:
        
        \begin{itemize}
            \item Quantum mechanics predictions
            \item Atomic spectroscopy
            \item Relativistic consistency
        \end{itemize}
    
    \subsection{Internal Consistency Rules}
        
        \begin{itemize}
            \item No circular definitions
            \item All primitives explicitly defined
            \item Derived quantities must follow from axioms
        \end{itemize}
    
    \subsection{Conflict Resolution Strategy}
        
        Priority order:
        
        \begin{enumerate}
            \item Logical consistency
            \item Experimental agreement
            \item Canonical simplicity
        \end{enumerate}
    
    \subsection{Epistemic Status Tracking}
        
        Each major result must explicitly state:
        
        \begin{itemize}
            \item Assumptions
            \item Derivation status
            \item Empirical testability
        \end{itemize}

        In particular, a definitional choice must not be presented as a theorem, and an interpretive bridge must not be presented as an orthodox result unless the underlying mathematical structure is itself standard.
    
    \subsection{Canon Integrity Principle}
        
        \begin{center}
            \textit{A valid SST canon must be internally consistent, externally compatible, and explicitly stratified by epistemic status.}
        \end{center}

    \section{Geometric and Topological Sector}
        \label{sec:geometry}

    \subsection{Role of geometry and topology in the canon}
        This sector collects those parts of SST in which the identity of a state depends not only on local field amplitudes but on the global organization of the carrier. Geometry enters through lengths, radii, and closure constraints; topology enters through linking, crossing, chirality, and other invariant features that survive smooth deformations. At canon level, the main role of this sector is selective rather than exhaustive: it provides the organizing variables that feed the topological kernel $\Xi_K$ and the geometric gate $\Pi_K$ in Eq.~\eqref{eq:sst_master_equation}.

        The canon therefore treats geometry and topology as follows:
        \begin{enumerate}
            \item geometry sets admissible scales and shielding ratios,
            \item topology labels persistent identity classes,
            \item only those quantities that survive smooth deformation are allowed to enter the long-lived state classification.
        \end{enumerate}

        This is also why layer or dressing indices are not treated as primary identity labels: they may modify a state, but they do not replace the topology-first classification of the carrier itself.

\subsection{Spectral Structure of Locked Closed Carriers (Secondary Result)}
    \label{subsec:spectral_locked_carriers}
    
    \paragraph{Scope and status.}
        This subsection records a \emph{secondary consequence} of the kinematics of rationally locked, closed carriers. It does not introduce new primitives. All results are conditional on the existence of a well-defined phase field and a regular mode spectrum. Status: \textbf{derived under spectral assumptions; experimentally unverified}.

    \paragraph{Kinematic reduction.}
        Consider a closed carrier with a single central driving phase $\phi(t)$ and $N_s$ sector variables $\theta_i(t)$ $(i=1,\dots,N_s)$ subject to a rational locking relation
        \begin{equation}
            \theta_i = \sigma_i \kappa_i \phi + \delta_i,
            \qquad
            \sigma_i \in \{+1,-1\},\quad \kappa_i \in \mathbb{Q}.
            \label{eq:locking_relation}
        \end{equation}
        In the rigid-lock limit, the dynamics reduce to a single collective coordinate $\phi(t)$ with effective Lagrangian
        \begin{equation}
            \mathcal{L}_{\mathrm{eff}}
            =
            \frac{I_{\mathrm{eff}}}{2}\dot{\phi}^{\,2}
            -
            V(\phi),
            \label{eq:Leff_canonical}
        \end{equation}
        where $I_{\mathrm{eff}}$ is an effective inertia and $V(\phi)$ is a periodic locking potential.

    \paragraph{Distributed phase field.}
        On a closed carrier of length $L$, allow small phase fluctuations
        \begin{equation}
            \phi(t) \;\longrightarrow\; \phi_0 + \eta(s,t),
            \qquad s\in[0,L],
        \end{equation}
        with periodic boundary condition
        \begin{equation}
            \eta(s+L,t)=\eta(s,t).
        \end{equation}
        A minimal quadratic field model consistent with Eq.~\eqref{eq:Leff_canonical} is
        \begin{equation}
            \mathcal{L}_{\mathrm{field}}
            =
            \int_0^L
            \left[
                \frac{\mu_{\mathrm{eff}}}{2}(\partial_t \eta)^2
            -
                \frac{T_{\mathrm{eff}}}{2}(\partial_s \eta)^2
            -
                \frac{\mu_{\mathrm{eff}}\omega_0^2}{2}\eta^2
            \right] ds,
            \label{eq:field_model}
        \end{equation}
        where $\mu_{\mathrm{eff}}$ is an effective phase inertia density, $T_{\mathrm{eff}}$ an effective stiffness, and $\omega_0$ the small-oscillation frequency set by $V(\phi)$.
        
        Dimensional consistency:
        \[
            [\mu_{\mathrm{eff}}]=\mathrm{kg\,m},\qquad
            [T_{\mathrm{eff}}]=\mathrm{N},\qquad
            [\omega_0]=\mathrm{s^{-1}}.
        \]

    \paragraph{Mode spectrum.}
        Expanding in normal modes
        \begin{equation}
            \eta(s,t) = \sum_{n\in\mathbb{Z}} a_n(t)\, e^{i k_n s},
            \qquad
            k_n = \frac{2\pi n}{L},
        \end{equation}
        yields the dispersion relation
        \begin{equation}
            \boxed{
                \omega_n^2
                =
                \omega_0^2
                +
                c_\phi^2 k_n^2,
                \qquad
                c_\phi^2 = \frac{T_{\mathrm{eff}}}{\mu_{\mathrm{eff}}}.
            }
            \label{eq:dispersion}
        \end{equation}
        For large $n$,
        \begin{equation}
            \omega_n \sim \frac{2\pi c_\phi}{L}\, n.
            \label{eq:asymptotic_linear}
        \end{equation}

    \paragraph{Cubic spectral sums and \texorpdfstring{$\zeta(3)$}{zeta(3)}.}
        Consider a cubic-weighted spectral contribution of the form
        \begin{equation}
            \Delta \mathcal{E}^{(3)}
            =
            A_3 \sum_{n=1}^{\infty} \frac{1}{\omega_n^3},
            \qquad
            [A_3]=\mathrm{J\,s^{-3}}.
            \label{eq:cubic_sum}
        \end{equation}
        Using the asymptotic scaling \eqref{eq:asymptotic_linear},
        \begin{equation}
            \Delta \mathcal{E}^{(3)}
            \;\approx\;
            A_3
            \left(\frac{L}{2\pi c_\phi}\right)^3
            \sum_{n=1}^{\infty} \frac{1}{n^3}.
        \end{equation}
        Hence
        \begin{equation}
            \boxed{
                \Delta \mathcal{E}^{(3)}
                \;\approx\;
                A_3
                \left(\frac{L}{2\pi c_\phi}\right)^3
                \zeta(3).
            }
            \label{eq:zeta3_result}
        \end{equation}

    \paragraph{Interpretation.}
        Equation \eqref{eq:zeta3_result} shows that $\zeta(3)$ arises from the \emph{high-mode spectral tail} of a closed carrier, not from the geometric threefold structure alone. The role of the locking relation \eqref{eq:locking_relation} is indirect: it fixes $I_{\mathrm{eff}}$, $\omega_0$, and admissible boundary conditions, thereby determining the spectral scale $(L,c_\phi)$.

    \paragraph{Universality.}
        Different geometric realizations (e.g., distinct multi-sector carriers) that reduce to the same class of effective phase dynamics and share the asymptotic linear spectrum \eqref{eq:asymptotic_linear} will produce the same $\zeta(3)$ factor in Eq.~\eqref{eq:zeta3_result}, up to system-dependent prefactors.
        
        \begin{equation}
            \boxed{
                \text{$\zeta(3)$ is a spectral invariant of the mode tower, not a direct invariant of the carrier geometry.}
            }
        \end{equation}

    \paragraph{Limitations and falsifiers.}
        The result fails if:
        \begin{enumerate}
            \item the mode spectrum is not asymptotically linear in $n$,
            \item strong damping or nonlocal coupling invalidates the expansion \eqref{eq:field_model},
            \item the system does not admit a well-defined periodic phase field on a closed domain.
        \end{enumerate}

    \paragraph{Minimal experimental test.}
        Measure mode frequencies $\omega_n$ for $n=1,2,3,\dots$ and verify
        \begin{equation}
            \frac{\omega_n}{n} \to \text{constant}
            \quad\text{as } n\to\infty.
        \end{equation}
        Deviation from this scaling excludes Eq.~\eqref{eq:zeta3_result}.

    \paragraph{Status.}
        \textbf{Secondary derived result.} Valid under spectral assumptions; not a defining postulate.

\subsection{Symmetry-Sensitive Dark-Knot Classification}
    \label{subsec:dark_knot_classification}

    A compact classification retained from earlier SST topology layers distinguishes visible, quasi-dark, and dark sectors by chirality, helicity, and low-order instability content.

    Let $\chi(K)$ denote a chirality indicator, let
    \begin{align}
        \mathcal{H}[K] = \int \mathbf{v}\cdot(\nabla\times\mathbf{v})\,dV
    \end{align}
    be the helicity functional, and let $N_u^{(1)}(K)$ count low-order unstable modes in a chosen perturbative stability analysis. Then the retained taxonomy is

    \begin{align}
        \text{dark:}&\quad \chi(K)=0,\quad \mathcal{H}[K]\approx 0,\quad N_u^{(1)}(K)=0, \\
        \text{quasi-dark:}&\quad |\chi(K)|\ll 1,\quad \mathcal{H}[K]\approx 0,\quad 0 < N_u^{(1)}(K)\le 2, \\
        \text{visible:}&\quad \chi(K)\neq 0 \quad \text{or}\quad \mathcal{H}[K]\neq 0.
        \label{eq:dark_knot_taxonomy}
    \end{align}

    \textbf{[ORTHODOX]}: helicity is a standard invariant of ideal vortex dynamics \cite{Moffatt1969,Batchelor1967}.

    \textbf{[SPECULATIVE]}: the interpretation of Eq.~\eqref{eq:dark_knot_taxonomy} as a matter-sector taxonomy is an SST model choice.

    \paragraph{Heuristic balance diagnostic from the SST helicity archive.}
    As a supporting diagnostic (not a stand-alone classifier), define
    \begin{align}
        \Delta_H(K)=\abs{\widehat{\mathcal{H}}_{b}(K)+\frac{1}{2}},
    \end{align}
    where $\widehat{\mathcal{H}}_{b}(K)$ is the archived normalized helicity-like base indicator for class $K$.
    Near-balanced values with $\Delta_H(K)\ll 1$ are treated only as candidate evidence for dark or quasi-dark sectors.

    \textbf{[CRITICAL NOTE]}: in canon usage, the diagnostic above is subordinate to Eq.~\eqref{eq:dark_knot_taxonomy}; no single scalar is sufficient for dark-sector membership without the symmetry and instability filters.

    \textbf{[CANON WORKFLOW NOTE]}: if SST-labelled and parallel legacy-labelled helicity exports are numerically identical, only the SST-labelled archive is required for canon workflows.

    \paragraph{Symmetry-first refinement (v0.8.x).}
        In the taxonomy program, candidate sectors are labeled symmetry-first, then filtered by helicity balance and low-order instability content.
        Amphichiral knot types furnish the preferred \emph{candidate} pool for the dark bin in symmetry-first tables; assigning ``dark'' still requires the conjunction of near-zero chirality indicators, helicity balance in the sense of Eq.~\eqref{eq:dark_knot_taxonomy}, and vanishing low-order instability count.
        The implication ``amphichiral $\Rightarrow$ dark'' is a classifier proposal tested against observables and archive diagnostics, not a closed theorem.

    \paragraph{Ambiguity tie-break.}
        If rule conjunctions overlap, precedence is
        \[
            \text{symmetry gate} \;\to\; \text{helicity balance gate} \;\to\; \text{stability gate}.
        \]

\section{Swirl--Electromagnetic Bridge}
    \label{sec:embridge}

    \subsection{Minimal transverse field sector}

        A conservative bridge from the fluid sector to a radiation sector is obtained by introducing a divergence-free transverse field $\mathbf{A}_{\mathrm{eff}}$ with

        \begin{align}
            \nabla\cdot \mathbf{A}_{\mathrm{eff}} = 0,
            \label{eq:div_free_Aeff}
        \end{align}

        and effective Lagrangian density

        \begin{align}
            \mathcal{L}_{\mathrm{rad}}
            =
            \frac{\rhoF}{2}\,\lVert \partial_t \mathbf{A}_{\mathrm{eff}}\rVert^2
            -
            \frac{K_\gamma}{2}\,\lVert \nabla\times\mathbf{A}_{\mathrm{eff}}\rVert^2.
            \label{eq:rad_lagrangian}
        \end{align}

        Varying Eq.~\eqref{eq:rad_lagrangian} and using Eq.~\eqref{eq:div_free_Aeff} gives

        \begin{align}
            \frac{1}{c_\gamma^2}\partial_t^2 \mathbf{A}_{\mathrm{eff}} - \nabla^2 \mathbf{A}_{\mathrm{eff}} = 0,
            \qquad
            c_\gamma^2 = \frac{K_\gamma}{\rhoF}.
            \label{eq:vector_wave_Aeff}
        \end{align}

        Define

        \begin{align}
            \mathbf{E}_{\mathrm{eff}} = -\partial_t \mathbf{A}_{\mathrm{eff}},
            \qquad
            \mathbf{B}_{\mathrm{eff}} = \nabla\times \mathbf{A}_{\mathrm{eff}}.
            \label{eq:effective_em_fields}
        \end{align}

        Then one immediately recovers the source-free identities

        \begin{align}
            \nabla\cdot \mathbf{B}_{\mathrm{eff}} &= 0, \\
            \nabla\times \mathbf{E}_{\mathrm{eff}} &= -\partial_t \mathbf{B}_{\mathrm{eff}}.
            \label{eq:maxwell_like_identities}
        \end{align}

        \textbf{[ORTHODOX FORM]}: Eq.~\eqref{eq:vector_wave_Aeff} and Eq.~\eqref{eq:maxwell_like_identities} are the standard transverse-wave and source-free Faraday identities in vector-potential form \cite{Jackson1999}.

        \textbf{[SPECULATIVE IDENTIFICATION]}: SST identifies $\mathbf{A}_{\mathrm{eff}}$ with a coarse-grained torsional or director displacement field of the medium.

    \subsection{Photon / unknot sector}

        The minimal retained radiation picture is that an R-phase excitation is an unknotted, finite-duration torsional packet propagating in the field $\mathbf{A}_{\mathrm{eff}}$. At the Rosetta level this corresponds to a source-free transverse wave packet with dispersion inherited from Eq.~\eqref{eq:vector_wave_Aeff}.

        \textbf{[SPECULATIVE]}: handedness of the local torsional packet is taken to encode polarization, while additional azimuthal phase winding encodes orbital angular momentum.

    \subsection{Swirl pressure law}
        \label{subsec:swirl_pressure_law}

        For steady axisymmetric swirl, Euler radial balance gives

        \begin{align}
            \frac{1}{\rhoF}\frac{dp_{\mathrm{swirl}}}{dr} = \frac{v_\theta(r)^2}{r}.
            \label{eq:swirl_pressure_law}
        \end{align}

        For a locally rigid rotation profile $v_\theta(r)=\Omega r$, Eq.~\eqref{eq:swirl_pressure_law} integrates to

        \begin{align}
            p_{\mathrm{swirl}}(r) = p_0 + \frac{1}{2}\rhoF \Omega^2 r^2,
            \label{eq:swirl_pressure_rigid}
        \end{align}

        while for an exterior irrotational profile $v_\theta(r)=\Gamma/(2\pi r)$ one finds

        \begin{align}
            p_{\mathrm{swirl}}(r) = p_\infty - \frac{\rhoF\Gamma^2}{8\pi^2 r^2}.
            \label{eq:swirl_pressure_irrotational}
        \end{align}

        In a flat rotation curve with asymptotically constant azimuthal speed $v_\theta(r)\to v_0$ as $r$ grows in a window where the radial balance \eqref{eq:swirl_pressure_law} remains valid, integration yields the logarithmic tail
        \begin{align}
            p_{\mathrm{swirl}}(r)
            =
            p_0 + \rhoF v_0^2 \ln\!\left(\frac{r}{r_0}\right),
            \label{eq:swirl_pressure_flat_tail}
        \end{align}
        for suitable reference radius $r_0$ and offset $p_0$ within that window.

        \textbf{[ORTHODOX] Assumption block:} Eq.~\eqref{eq:swirl_pressure_law} is taken in the usual incompressible, inviscid, stationary, axisymmetric regime with azimuthal-dominant flow, so that radial drift, viscosity, and stratification corrections are subleading in the domain of validity.

        \textbf{[ORTHODOX] Dimensional check:}
        $\bigl[\rhoF^{-1}\,dp/dr\bigr]=\mathrm{m\,s^{-2}}=\bigl[v_\theta^2/r\bigr]$.

        \textbf{[ORTHODOX]}: Eq.~\eqref{eq:swirl_pressure_law} is the ordinary radial Euler balance for steady swirl \cite{Batchelor1967,Saffman1992}.

        \textbf{[DERIVED SST ROLE]}: this pressure law is the canonical bridge between localized circulation and large-scale force analogies. Interpreting $p_{\mathrm{swirl}}$ as a dark-sector classifier or support layer is an SST model step, logically separate from the orthodox fluid derivation.

        \textbf{[DERIVED] Observables / falsifiers:} (i) consistency of a pressure profile reconstructed from archived $v_\theta(r)$ with the radial gradient implied by \eqref{eq:swirl_pressure_law}; (ii) null test: systematic sign or scale mismatch of the reconstructed $dp_{\mathrm{swirl}}/dr$ relative to $v_\theta^2/r$ across validated radial windows excludes the applicability of the balance in those windows.

    \subsection{Minimal effective Lagrangian bridge}

        A compact field-theoretic summary retaining the clock, radiation, and matter sectors is

        \begin{align}
            \mathcal{L}_{\mathrm{eff}}
            =
            \mathcal{L}_{T}
            +
            \mathcal{L}_{\mathrm{rad}}
            +
            \sum_K \bar\Psi_K\bigl(i\gamma^\mu D_\mu - m_K^{(\mathrm{sol})}\bigr)\Psi_K
            + \cdots,
            \label{eq:minimal_effective_lagrangian_bridge}
        \end{align}

        where $\mathcal{L}_{T}$ is the clock/foliation sector introduced in Section~\ref{sec:time}, $\mathcal{L}_{\mathrm{rad}}$ is given by Eq.~\eqref{eq:rad_lagrangian}, and the ellipsis denotes higher-order self-interaction, topological, or medium-response operators.

        \textbf{[SPECULATIVE BRIDGE]}: Eq.~\eqref{eq:minimal_effective_lagrangian_bridge} is a bookkeeping layer, not yet a completed first-principles derivation.

% ---------- DELAY THEORY ----------
%    \section{Delay and Mode Selection Theory}
%        \label{sec:delay}
        
\section{Delay and Mode Selection Theory}
    \label{sec:delay}
    
    \subsection{Circulation Delay as a Physical Primitive}
        
        We introduce the circulation delay as a constitutive element of the theory.
        For a closed vortex filament of effective length $L$ and characteristic propagation speed $\vnorm$, the circulation time is defined as:
        
        \begin{align}
            \tau = \frac{L}{\vnorm}
            \label{eq:delay_time}
        \end{align}
        
        In the special case of a minimal closed loop associated with the vortex core scale, this reduces to:
        
        \begin{align}
            \tau_c = \frac{2\pi \rc}{\vnorm} = \frac{2\pi}{\omega_c}
            \label{eq:compton_delay}
        \end{align}
        
        This establishes a direct identification between the circulation delay and the inverse Compton frequency.
        
        \textbf{[DERIVED]}: The intrinsic delay of the vortex loop is identical to the Compton period.
    
    \subsection{Delayed Self-Interaction Dynamics}
        
        We model the phase $\phi(t)$ of a circulating excitation as a delayed self-interacting oscillator:
        
        \begin{align}
            \dot{\phi}(t) = \omega_0 + \kappa \sin\big(\phi(t - \tau) - \phi(t)\big)
            \label{eq:delay_phase}
        \end{align}
        
        This equation represents a minimal effective description of a circulating field with finite propagation time.
        
        \textbf{[ORTHODOX]}: Delay-differential equations of this type are well-established in feedback systems \cite{SST_delay_basic}.
    
    \subsection{Emergence of Discrete Stable Modes}
        
        We seek steady-state solutions of the form:
        
        \begin{align}
            \phi(t) = \Omega t
        \end{align}
        
        Substituting into Eq.~\ref{eq:delay_phase} yields:
        
        \begin{align}
            \Omega = \omega_0 + \kappa \sin(\Omega \tau)
            \label{eq:mode_condition}
        \end{align}
        
        This transcendental equation admits a discrete set of stable solutions $\Omega_n$ determined by the delay parameter $\tau$.
        
        \textbf{[DERIVED]}: Discrete mode selection arises from temporal nonlocality alone.
    
    \subsection{Delay-Induced Quantization Mechanism}
        
        Unlike conventional quantum mechanics, where discreteness is postulated, here it emerges dynamically:
        
        \begin{itemize}
            \item No boundary condition required
            \item No operator quantization assumed
            \item Discreteness arises from stability selection
        \end{itemize}
        
        \textbf{[SPECULATIVE]}: Particle spectra may be interpreted as delay-selected stable circulation modes.
    
    \subsection{Extension to Nonlinear Schrödinger Dynamics}
        
        To connect with field-theoretic descriptions, we extend the delay mechanism to a nonlinear Schrödinger framework:
        
        \begin{align}
            i \hbar \partial_t \psi = -\frac{\hbar^2}{2m} \nabla^2 \psi + F_{\text{delay}}[\psi(t - \tau)]
        \end{align}
        
        where the delayed functional introduces temporal nonlocality into the field evolution.
        
        \textbf{[SPECULATIVE]}: This provides a bridge between classical delay systems and quantum wave dynamics.
    
    \subsection{Physical Interpretation}
        
        The delay $\tau$ acts as:
        
        \begin{itemize}
            \item a memory kernel
            \item a nonlocal temporal coupling
            \item a selection mechanism for stable modes
        \end{itemize}
        
        Thus:
        
        \begin{center}
            \textit{Delay $\Rightarrow$ Stability $\Rightarrow$ Discreteness}
        \end{center}
        
        This mechanism forms one of the central pillars of the SST canon.

%    \section{Atomic Bridge Model}
%        \label{sec:atomic}
        
\section{Atomic Bridge Model}
    \label{sec:atomic}

    \subsection{Introduction and epistemic framing}
        This section records the minimal SST atomic bridge rather than a full derivation of atomic structure. Its canon role is deliberately narrow: recover the orthodox Coulomb limit when the SST-specific sector is switched off, define a branch-resolved circulation coupling fixed by canonical medium scales, and make explicit which many-electron and shell-structure claims remain open. It is therefore a benchmark layer, not yet a complete atomic theory.
    
    \subsection{Baseline Hamiltonian}
        
        We adopt the standard Coulomb Hamiltonian as the orthodox baseline:
        
        \begin{align}
            H_0 = -\frac{\hbar^2}{2\mu} \nabla^2 - \frac{Ze^2}{4\pi \epsilon_0 r}
            \label{eq:coulomb}
        \end{align}
        
        \textbf{[ORTHODOX]}: This reproduces the hydrogenic spectrum.

    \subsection{Legacy Soft-Core Regularization}

        Early canon layers introduced a finite-core regularization of the hydrogenic potential. In the present document this is retained only as a controlled bridge ansatz, not as a replacement for Eq.~\eqref{eq:coulomb}.

        Define the effective coupling scale
        \begin{align}
            \Lambda_{\mathrm{SST}} = 4\pi \, \rhocore \, \vnorm^{2} \, \rc^{4}
            \label{eq:lambda_sst}
        \end{align}
        with dimensions
        \begin{align}
            [\Lambda_{\mathrm{SST}}] = \mathrm{J\,m} = \mathrm{N\,m^2}.
        \end{align}

        A finite-core bridge potential is then
        \begin{align}
            V_{\mathrm{soft}}(r) = -\frac{\Lambda_{\mathrm{SST}}}{\sqrt{r^2+r_c^2}}
            \label{eq:softcore_potential}
        \end{align}
        satisfying the large-radius limit
        \begin{align}
            V_{\mathrm{soft}}(r) \xrightarrow{r\gg r_c} -\frac{\Lambda_{\mathrm{SST}}}{r}.
        \end{align}

        \textbf{[ORTHODOX]} Effective-model input: finite-core regularizations of singular central potentials are standard tools.\newline
        \textbf{[DERIVED]} Canonical consequence: the SST scale in Eq.~\eqref{eq:lambda_sst} is fixed by canonical constants rather than introduced phenomenologically.\newline
        \textbf{[ORTHODOX] Constraint}: this sector is admissible only if the large-radius limit reproduces the standard Coulomb problem to spectroscopic accuracy.

        In particular, Eq.~\eqref{eq:softcore_potential} is retained as a bridge ansatz only. It is not promoted here to the primary atomic potential.

    \subsection{SST Circulation Correction}
        
        We introduce a circulation-dependent perturbation:
        
        \begin{align}
            \delta H^{(\sigma)} = -2\sigma E_R \eta_0 \gamma + \frac{E_R \eta_0 \sigma}{2} \left\{ \frac{\hat{L}_z}{\hbar}, \gamma \right\}
            \label{eq:sst_correction}
        \end{align}
        
        \textbf{[DERIVED]}: Interaction depends on circulation amplitude $\gamma(r,t)$.

    \subsection{Circulation source chain}
        The atomic perturbation is not introduced as an arbitrary extra scalar, but is intended to descend from a circulation-bearing source chain,
        \begin{align}
            \xi \rightarrow \omega_z \rightarrow u_\theta \rightarrow \Gamma_{\mathrm{loc}} \rightarrow \gamma,
        \end{align}
        where an axial displacement field $\xi$ generates local vorticity, the vorticity sources azimuthal swirl, and the loop circulation is normalized by the canonical quantum $\Gamma_0$. This makes the circulation amplitude $\gamma$ the natural canon-level coupling variable rather than a free phenomenological insertion.
    
    \subsection{Branch Structure}
        
        Two chiral sectors exist:
        
        \begin{align}
            \sigma = \pm 1
        \end{align}
        
        These correspond to:
        
        \begin{itemize}
            \item $\swirlarrow$ (left-handed)
            \item $\swirlarrowcw$ (right-handed)
        \end{itemize}
    
    \subsection{Consistency Requirement}
        
        To remain compatible with atomic physics:
        
        \begin{align}
            |\delta E_n| \ll |E_n|
        \end{align}
        
        \textbf{[ORTHODOX] Constraint}: SST corrections must not violate known spectroscopy.
    
    \subsection{Interpretation}
        
        \textbf{[STATUS]}: The atomic bridge is a controlled benchmark sector. The orthodox Coulomb problem is imported as the baseline; the circulation-dependent term is the SST addition; shell counting, exclusion-like behavior, and heavy-atom extensions remain programmatic rather than derived here. The purpose of the section is therefore structural compatibility, not yet a complete replacement of atomic quantum mechanics.

%    \section{Spectroscopic Constraints}
%        \label{sec:spectroscopy}
        
\section{Spectroscopic Constraints}
    \label{sec:spectroscopy}
    
    \subsection{Motivation}
        
        Any extension of atomic structure must remain consistent with high-precision spectroscopic measurements.
        
        \textbf{[ORTHODOX]}:
        Atomic energy levels are experimentally verified to extremely high precision.
        
        \textbf{Requirement:}
        \begin{align}
            |\delta E_n| \ll |E_n|
            \label{eq:spectro_constraint}
        \end{align}
        
        where $\delta E_n$ represents any SST-induced correction.
    
    \subsection{Internal Mode Hypothesis}
        
        We consider the possibility that atomic bound states possess an additional internal excitation sector associated with vortex filament dynamics.
        
        \textbf{[SPECULATIVE]}:
        Each bound state carries an internal mode spectrum $\{E_{m,n}\}$.
        
        The effective energy becomes:
        
        \begin{align}
            E_n^{\text{eff}} = E_n^{(0)} + \delta E_n
            \label{eq:effective_energy}
        \end{align}
    
    \subsection{Generic Coupling Structure}
        
        We parameterize the coupling between internal modes and observable energy levels:
        
        \begin{align}
            |\delta E_n| \leq \lambda_n \sum_m E_{m,n} \, \nu_{m,n}
            \label{eq:coupling_bound}
        \end{align}
        
        where:
        \begin{itemize}
            \item $\lambda_n$ is the coupling efficiency
            \item $\nu_{m,n}$ is the occupation factor
        \end{itemize}
        
        \textbf{[DERIVED]}: This form captures a general upper bound independent of microscopic details.
    
    \subsection{Ungapped Spectrum: Instability}
        
        If the internal spectrum is ungapped:
        
        \begin{itemize}
            \item arbitrarily low-energy excitations exist
            \item no universal suppression mechanism applies
        \end{itemize}
        
        \textbf{Consequence:}
        \begin{align}
            \delta E_n \not\to 0 \quad \text{generically}
        \end{align}
        
        \textbf{[CRITICAL]}:
        An ungapped internal sector is generically incompatible with spectroscopy unless:
        
        \begin{itemize}
            \item $\lambda_n \ll 1$ (extremely weak coupling), or
            \item the density of states is highly suppressed
        \end{itemize}
    
    \subsection{Gapped Spectrum: Decoupling Mechanism}
        
        We introduce an excitation gap $\Delta_K$:
        
        \begin{align}
            E_{m,n} \geq \Delta_K
        \end{align}
        
        Then occupation factors follow Boltzmann suppression:
        
        \begin{align}
            \nu_{m,n} \sim \exp\left(-\frac{\Delta_K}{k_B T_{\text{eff}}}\right)
        \end{align}
        
        Substituting into Eq.~\ref{eq:coupling_bound}:
        
        \begin{align}
            \delta E_n \propto \exp\left(-\frac{\Delta_K}{k_B T_{\text{eff}}}\right)
            \label{eq:boltzmann_suppression}
        \end{align}
        
        \textbf{[DERIVED]}:
        Exponential suppression ensures compatibility with spectroscopy.
    
    \subsection{Physical Origin of the Gap}
        
        \textbf{[SPECULATIVE]}:
        The gap $\Delta_K$ may arise from:
        
        \begin{itemize}
            \item geometric constraints on filament curvature
            \item torsional rigidity
            \item topological restrictions
            \item core structure of the vortex
        \end{itemize}
    
    \subsection{Order-of-Magnitude Estimate}
        
        Using filament parameters:
        
        \begin{align}
            \Delta_K \sim \frac{\Gamma^2}{4\pi \xi L}
        \end{align}
        
        where:
        \begin{itemize}
            \item $\xi$ is the core scale
            \item $L$ is the loop length
        \end{itemize}
        
        For typical SST scales:
        
        \begin{align}
            \Delta_K = \mathcal{O}(10^2 - 10^3 \, \text{eV})
        \end{align}
        
        \textbf{[DERIVED]}:
        This places the internal sector well above atomic energy scales ($\sim$ eV).
    
    \subsection{Consistency Condition}
        
        To preserve atomic physics:
        
        \begin{align}
            \frac{\Delta_K}{k_B T_{\text{eff}}} \gg 1
        \end{align}
        
        \textbf{Interpretation:}
        The internal sector is effectively frozen at atomic energies.
    
    \subsection{Observable Consequences}
        
        \textbf{[SPECULATIVE]}:
        Residual effects may include:
        
        \begin{itemize}
            \item tiny level shifts
            \item suppressed line broadening
            \item weak temperature dependence
        \end{itemize}
    
    \subsection{Falsifiability}
        
        The theory is falsifiable via:
        
        \begin{itemize}
            \item detection of unexplained spectral shifts
            \item absence of expected suppression
            \item deviation from Boltzmann scaling
        \end{itemize}
    
    \subsection{Integration with SST Framework}
        
        The spectroscopic constraint enforces:
        
        \begin{itemize}
            \item compatibility with the Atomic Bridge (Section~\ref{sec:atomic})
            \item constraints on delay-induced excitations (Section~\ref{sec:delay})
            \item bounds on allowable vortex configurations
        \end{itemize}
    
    \subsection{Canonical Statement}
        
        \begin{center}
            \textit{
                Any viable SST extension must include a mechanism—such as an excitation gap—that ensures exponential suppression of internal modes at atomic energy scales.
            }
        \end{center}

    \subsection{Precision electroweak bridge: lessons from the CMS \texorpdfstring{$W$}{W}-mass measurement}
        \label{subsec:cms-w-mass-bridge}
    
        \paragraph{Status.}
            \textbf{[ORTHODOX]} External input: the CMS measurement is a precision constraint on any SST effective weak-sector realization.
            \textbf{[DERIVED]} Methodological consequence: the result supports a forward-modelling strategy for SST phenomenology.
            \textbf{[CRITICAL NOTE]} The paper does \emph{not} provide direct evidence for a structured-substrate ontology of the $W$ boson.
    
        \paragraph{Empirical anchor.}
            Let the experimentally inferred resonance mass be denoted by
            \begin{equation}
                m_W^{\rm exp} = 80.3602 \pm 0.0099~{\rm GeV}.
            \end{equation}
            In the CMS analysis this quantity is not read off from a single observable, but extracted from a profiled likelihood over finely binned charged-lepton kinematics. This motivates the following SST principle:
            
            \begin{quote}
                \emph{Collider observables constrain SST through forward maps from latent swirl-string states to detector-level distributions, not through direct identification of a single internal parameter with a measured mass.}
            \end{quote}
    
        \paragraph{Latent-to-observable map.}
            Introduce a latent SST parameter set for an effective spin-1 weak excitation,
            \begin{equation}
                \Theta_{W}^{\rm SST}
                =
                \Bigl(
                K_W,\,
                \chi_W,\,
                \mathcal{L}_W,\,
                \mathcal{H}_W,\,
                \SwirlClock_W,\,
                \rhoF^{\rm eff},\,
                \vswirl^{\rm eff}
                \Bigr),
            \end{equation}
            where:
            \begin{itemize}
                \item $K_W$ is the topology class or excitation family,
                \item $\chi_W$ is chirality / handedness label,
                \item $\mathcal{L}_W$ is an effective geometric length scale,
                \item $\mathcal{H}_W$ is a helicity/linkage content,
                \item $\SwirlClock_W$ is the local clock factor for the mode,
                \item $\rhoF^{\rm eff}$ and $\vswirl^{\rm eff}$ are the coarse-grained medium parameters seen by the excitation.
            \end{itemize}
            
            The experimentally accessible distribution is then not a bare function of $m_W$ alone, but a projected density
            \begin{equation}
                \mathcal{P}_{\rm meas}(x)
                =
                \int dz\; R(x|z)\,\mathcal{P}_{\rm prod}(z\mid \Theta_W^{\rm SST}),
                \label{eq:sst-forward-map}
            \end{equation}
            where $z$ denotes truth-level production/decay variables and $R(x|z)$ is the detector response kernel.
    
        \paragraph{Mass as an effective resonance parameter.}
            In SST one should distinguish:
            \begin{equation}
                m_W^{\rm internal}
                \neq
                m_W^{\rm eff}
            \end{equation}
            in general, where:
            \begin{itemize}
                \item $m_W^{\rm internal}$ is the internal excitation-energy scale of the structured mode,
                \item $m_W^{\rm eff}$ is the resonance parameter inferred from scattering distributions.
            \end{itemize}
            The experimentally constrained quantity is $m_W^{\rm eff}$.
            
            A minimal canon-safe matching condition is therefore
            \begin{equation}
                m_W^{\rm eff}(\Theta_W^{\rm SST}) = m_W^{\rm exp} + \delta_W,
                \qquad
                |\delta_W| \lesssim 10~{\rm MeV},
                \label{eq:mw-match}
            \end{equation}
            with the current precision scale set by experiment.
    
        \paragraph{Distribution-level fit target.}
            The correct SST comparison object is a likelihood built from binned observables:
            \begin{equation}
                \mathcal{L}\bigl(\Theta_W^{\rm SST},\nu\bigr)
                =
                \prod_{b=1}^{N_{\rm bins}}
                    {\rm Poisson}\!\left(
                                       n_b \,\middle|\, \mu_b(\Theta_W^{\rm SST},\nu)
                \right)
                \,
                \Pi(\nu),
                \label{eq:sst-likelihood}
            \end{equation}
            where:
            \begin{itemize}
                \item $n_b$ is the observed count in bin $b$,
                \item $\mu_b$ is the SST prediction after production, decay, and detector folding,
                \item $\nu$ denotes nuisance parameters,
                \item $\Pi(\nu)$ is the nuisance prior / constraint factor.
            \end{itemize}
            The profile likelihood for the physically interesting SST parameters is then
            \begin{equation}
                \mathcal{L}_{\rm prof}(\Theta_W^{\rm SST})
                =
                \max_{\nu}\,
                \mathcal{L}(\Theta_W^{\rm SST},\nu).
                \label{eq:profile-like}
            \end{equation}
    
        \paragraph{Helicity decomposition as a falsifier.}
            A particularly sharp SST test is not the central value of $m_W$ alone, but whether the effective spin-1 mode reproduces the standard angular basis. Write
            \begin{equation}
                \frac{d\sigma}{d\Omega}
                =
                \sum_{i=0}^{8} A_i\, f_i(\theta,\phi)
                +
                \epsilon_{\rm SST}\,\Delta(\theta,\phi),
                \label{eq:helicity-decomp-sst}
            \end{equation}
            where the $A_i f_i$ are the standard helicity structures and $\epsilon_{\rm SST}\Delta$ is the first SST correction. Then:
            \begin{itemize}
                \item if $\epsilon_{\rm SST}=0$ within experimental precision, SST survives this channel,
                \item if $\epsilon_{\rm SST}\neq 0$, the deviation is directly falsifiable through angular distributions.
            \end{itemize}
    
        \paragraph{Canon interpretation.}
            This leads to the following conservative reading:
            \begin{enumerate}
                \item The weak sector must be \emph{SM-like at the effective distribution level}.
                \item Large $\mathcal{O}(10^{-4})$--$\mathcal{O}(10^{-3})$ relative distortions in the effective $W$ mass are disfavoured.
                \item The best SST discovery channels are therefore likely to be
                \begin{equation}
                    \text{small angular deviations},\quad
                    \text{polarization asymmetries},\quad
                    \text{off-shell tails},\quad
                    \text{rare differential distortions},
                \end{equation}
                rather than a large shift in the central value of $m_W$.
            \end{enumerate}
    
        \paragraph{Minimal bridge claim.}
            The CMS result does \emph{not} force the $W$ boson to be fundamental in the ontological sense. It does force any SST realization of the weak sector to satisfy the stronger requirement:
            \begin{equation}
                \mathcal{P}_{\rm prod}^{\rm SST}(z\mid \Theta_W^{\rm SST})
                \approx
                \mathcal{P}_{\rm prod}^{\rm SM}(z\mid m_W,\text{SM nuisances})
            \end{equation}
            to current experimental precision after detector folding and nuisance profiling.
    
        \paragraph{Practical SST program.}
            A precision-compatible SST weak-sector program should therefore proceed in three layers:
            \begin{align}
                \text{(i) effective resonance matching:}&\quad
                m_W^{\rm eff} \to m_W^{\rm exp}, \\
                \text{(ii) helicity matching:}&\quad
                A_i^{\rm SST} \to A_i^{\rm SM}, \\
                \text{(iii) residual search:}&\quad
                \Delta \neq 0 \text{ only in subleading differential structure.}
            \end{align}
    
        \paragraph{Conclusion.}
            The CMS $W$-mass result is best understood in SST as a \emph{precision filter}:
            it does not supply ontology, but it strongly constrains phenomenology. Any viable SST weak-sector realization must be distributionally SM-like to present precision, while leaving room only for small, structured residuals in higher-order observables.

% ---------- TIME ----------
    \section{Relational Time Framework}
        \label{sec:time}

    \subsection{Swirl-Clock and local proper time}
        The canonical time variable of SST is not introduced as an independent primitive. It is defined relationally from the local kinematic state of the medium through the Swirl-Clock already introduced in Eq.~\eqref{eq:swirl_clock}. In its coarse-grained form,
        \begin{align}
            dt_{\mathrm{loc}} = S_{(t)}(x)\,dt_{\infty},
            \qquad
            S_{(t)}(x)=\sqrt{1-\frac{v(x)^2}{c^2}}.
        \end{align}
        Here $t_{\infty}$ denotes the asymptotic laboratory time and $t_{\mathrm{loc}}$ the effective proper time associated with a localized process in the medium. In SST, this factor is interpreted as a fluid-kinematic clock field rather than as a purely geometric spacetime postulate.

    \subsection{Carrier-local clocking and the turnover scale}
        For localized carriers the natural microscopic rate is the turnover frequency
        \begin{align}
            \Omega_0 = \frac{\vnorm}{\rc} = \omega_c,
        \end{align}
        which coincides with the Compton anchoring scale. A complete SST clock description therefore involves two linked but distinct levels:
        \begin{itemize}
            \item the \emph{microscopic turnover scale} $\Omega_0$, which counts circulation cycles of the carrier,
            \item the \emph{coarse-grained Swirl-Clock} $S_{(t)}$, which relates local rates to asymptotic time.
        \end{itemize}
        This resolves a recurring ambiguity in earlier canon versions: local cycle counting and coarse-grained time dilation are related, but they are not identical observables.

    \subsection{Relational observables from conserved currents}
        A purely operator-based notion of time is avoided in SST by defining time observables through conserved currents. Let $J^\mu$ denote a matter current and let $j_{\mathrm{ev}}^\mu$ denote a coarse-grained event current associated with structured activity of the medium. The relational clock field $T(x)$ is then recovered only after coarse graining,
        \begin{align}
            \partial_\mu T(x) \approx \frac{1}{\nu_0}\,\langle j^{\mathrm{ev}}_\mu \rangle_{\ell},
        \end{align}
        where $\ell$ is the coarse-graining scale and $\nu_0$ a reference event frequency. The canon interpretation is that time is read from correlations between currents, not from an abstract external operator.

    \subsection{Time-of-arrival as a relational field observable}
        Let $\mathcal{W}$ be a detector world-tube with boundary $\Sigma = \partial \mathcal{W}$. Then the arrival-time distribution associated with a detector label $\Theta$ is written as
        \begin{align}
            p(\Theta) = \frac{1}{\mathcal{N}} \int_{\Sigma} d\sigma_\mu \, J^\mu(x)\,\delta\!\big(T(x)-\Theta\big).
        \end{align}
        This formulation is canonically important because it bypasses the need for a self-adjoint time operator conjugate to a bounded Hamiltonian. The measured time is instead the value of a physical clock field at the point where matter flux crosses the detector boundary.

    \subsection{Clock broadening and continuum limits}
        If the event current has a finite variance, then the observed time-of-arrival distribution is broadened relative to the ideal classical prediction. Writing $P_{\mathrm{cl}}$ for the classical arrival distribution and $\sigma_\tau^2$ for the effective clock variance, one expects a convolution structure of the form
        \begin{align}
            P_{\mathrm{obs}}(\Theta)
            =
            \int dt\,P_{\mathrm{cl}}(t)
            \frac{1}{\sqrt{2\pi\sigma_\tau^2}}
            \exp\!\left[-\frac{(\Theta-t)^2}{2\sigma_\tau^2}\right].
        \end{align}
        In the limit $\ell \gg r_c$, the event structure is smoothed and one recovers a continuum clock field. In the opposite limit $\ell \to r_c$, the underlying discreteness of carrier events becomes operationally relevant.

    \subsection{Connection to gravity and geometry}
        In the canon, the relational-time sector and the geometric-response sector are not independent. Variations in the Swirl-Clock track variations in local swirl energy density and therefore provide the time component of the effective geometric response of the medium. At this level, the SST reading is:
        \begin{center}
            \textit{gravity and time are two coarse-grained aspects of the same clock-bearing medium.}
        \end{center}
        This statement should be read as a canon-level organizing principle, not yet as a completed field-theoretic proof.

    \subsection{Clock Field and Foliation Variable}

        \textbf{[DERIVED]} EFT bridge: the scalar Swirl-Clock field is promoted from the kinematic relation
        of Eq.~\eqref{eq:swirl_clock} to a long-wavelength foliation variable
        \begin{align}
            \chi(x),
        \end{align}
        whose normalized gradient defines a preferred local clock frame:
        \begin{align}
            u_\mu
            =
            \frac{\nabla_\mu \chi}
            {\sqrt{g^{\alpha\beta}\nabla_\alpha\chi\,\nabla_\beta\chi}}.
            \label{eq:clock_foliation_vector}
        \end{align}

        \textbf{[ORTHODOX]} Equation~\eqref{eq:clock_foliation_vector} matches the standard hypersurface-orthogonal
        foliation construction used in khronometric / Einstein--\AE ther effective field theory
        \cite{Jacobson2008,BlasPujolasSibiryakov2011}.

        \textbf{[DERIVED]} Within SST, the interpretation is different: \(u_\mu\) is not taken as a fundamental
        spacetime structure, but as the coarse-grained clock-flow direction associated with the medium.

    \subsection{Clock-Sector Effective Action}

        \textbf{[ORTHODOX]} EFT template: the most general second-order covariant clock-sector action compatible
        with hypersurface orthogonality is
        \begin{align}
        S_\chi
        =
        \int d^4x \sqrt{-g}\,
        \Big[
        & c_1 (\nabla_\mu u_\nu)(\nabla^\mu u^\nu)
        + c_2 (\nabla_\mu u^\mu)^2 \nonumber\\
        & + c_3 (\nabla_\mu u_\nu)(\nabla^\nu u^\mu)
        + c_4 u^\mu u^\nu (\nabla_\mu u_\alpha)(\nabla_\nu u^\alpha)
        \Big].
        \label{eq:clock_sector_action}
        \end{align}

        \textbf{[SPECULATIVE]} SST interpretation: Equation~\eqref{eq:clock_sector_action} is not itself derived
        from first-principles filament dynamics in the present Canon. It is recorded as the minimal orthodox EFT
        bridge into a relativistic clock-field description.

    \subsection{Poisson Limit and Effective Gravity Bridge}

        \textbf{[DERIVED]} Bridge statement: in the weak-field, slowly varying limit, the scalar clock field is taken to obey
        a Poisson-type closure
        \begin{align}
            \nabla^2 \chi(x)
            =
            4\pi G_{\mathrm{eff}} \,\rho_{\mathrm{matter}}(x),
            \label{eq:poisson_clock}
        \end{align}
        so that outside localized sources
        \begin{align}
            \chi(r) \propto \frac{1}{r},
            \qquad
            \mathbf{g}_{\mathrm{eff}} = - \nabla \chi.
            \label{eq:effective_gravity_clock}
        \end{align}

        \textbf{[ORTHODOX]} The \(1/r\) scalar potential and \(1/r^2\) force law are standard consequences of
        Eq.~\eqref{eq:poisson_clock}.

        \textbf{[SPECULATIVE]} The identification of \(\chi\) with the Swirl-Clock and the interpretation of
        Eq.~\eqref{eq:poisson_clock} as an emergent gravity law are SST-specific.

    \subsection{Observational Constraint from Tensor-Mode Speed}

        \textbf{[ORTHODOX] Constraint}: If the clock-sector EFT is used, multimessenger gravitational-wave
        observations require \cite{Monitor2017}
        \begin{align}
            c_{13} \equiv c_1 + c_3 = 0,
            \label{eq:c13_constraint}
        \end{align}
        to enforce luminal tensor-mode propagation at observational precision.

        \textbf{[ORTHODOX] Constraint}: any SST realization employing the EFT bridge of Eq.~\eqref{eq:clock_sector_action}
        must preserve Eq.~\eqref{eq:c13_constraint}.

    \subsection{Interpretive Role in the Canon}

        This section has a deliberately limited role:
        \begin{itemize}
            \item it upgrades the Swirl-Clock from a purely local kinematic factor to a candidate field variable;
            \item it records a minimal orthodox EFT bridge for relativistic comparison;
            \item it supplies the internal reference point needed by Section~\ref{sec:unification} for the time--gravity link.
        \end{itemize}

        \textbf{Status summary.}
        \begin{itemize}
            \item Eq.~\eqref{eq:clock_foliation_vector}: \textbf{[ORTHODOX]} mathematical form; \textbf{[DERIVED]} role as an SST bridge variable
            \item Eq.~\eqref{eq:clock_sector_action}: \textbf{[ORTHODOX]} EFT template
            \item Eq.~\eqref{eq:poisson_clock}: \textbf{[DERIVED]} bridge closure
            \item gravity interpretation: \textbf{[SPECULATIVE]}
        \end{itemize}

    \subsection{Relation to the Chronos--Kelvin law}

        Section~\ref{subsec:chronos_kelvin} provides the local transport identity of the Swirl-Clock on a material loop. The present section supplies the long-wavelength field-theoretic and analogue-metric reading of the same clock sector. The two layers are therefore complementary:

        \begin{itemize}
            \item Eq.~\eqref{eq:clock_radius_transport} is the loop-level transport law;
            \item Eq.~\eqref{eq:clock_foliation_vector}--Eq.~\eqref{eq:poisson_clock} provide the coarse-grained field description.
        \end{itemize}

%    \section{Unified Interpretation}
%        \label{sec:unification}
        
\section{Unified Interpretation}
    \label{sec:unification}
    
    \subsection{Overview}
        
        Swirl-String Theory (SST) proposes a unified interpretation in which fundamental physical phenomena emerge from the dynamics of a continuous medium supporting stable vortex structures.
        
        This section synthesizes the results of the previous sections into a coherent conceptual framework.
    
    \subsection{Ontology of the Theory}
        
        \textbf{[SPECULATIVE] Ontological claim}:
        The fundamental ontology consists of:
        
        \begin{itemize}
            \item A continuous incompressible medium (Section~\ref{sec:axioms})
            \item Stable vortex filament configurations (Section~\ref{sec:geometry})
            \item Circulation as a conserved topological invariant
        \end{itemize}
        
        \textbf{Interpretation:}
        Particles are not elementary objects, but persistent dynamical structures.
    
    \subsection{Mass as Trapped Circulation Energy}
        
        From Section~\ref{sec:foundations}, and in particular from Eq.~\eqref{eq:sst_master_equation}, the mass sector is organized as the product of a medium scale, a geometric gate, a topological kernel, and a local clock-impedance factor. The canon-level reading is therefore not merely that ``energy is stored in rotation,'' but that persistent inertial scale is assigned only after geometry, topology, and local clocking have all been specified.
        
        \begin{align}
            M_K(x)=\rhoM(x)\,\rc^3\,\Pi_K\,\Xi_K\,\SwirlClock(x)^{-2}.
        \end{align}
        
        \textbf{[SPECULATIVE] Unified interpretation:} mass corresponds to localized, self-sustained circulation energy in the medium, modulated by state geometry, topology, and local clock impedance.
    
    \subsection{Time as a Relational Quantity}
        
        From Section~\ref{sec:time} and Section~\ref{sec:foundations}, time is not introduced as an independent primitive but as a relational field read from the local state of the medium. The same kinematic quantity that enters the Swirl-Clock also appears in the organizing equation of the mass sector, so the canon treats proper-time modulation and inertial scaling as two projections of one shared substrate state.
        
        \begin{align}
            \SwirlClock(x)=\sqrt{1-\frac{\norm{\vswirl(x)}^2}{c^2}}.
        \end{align}
        
        \textbf{[SPECULATIVE] Unified interpretation:} time is not fundamental, but emerges from the kinematic state of the medium and is read operationally through clock-bearing processes.
    
    \subsection{Discreteness without Quantization Postulates}
        
        From Section~\ref{sec:delay}:
        
        \begin{itemize}
            \item Delay introduces temporal nonlocality
            \item Stability selects discrete modes
        \end{itemize}
        
        \textbf{[DERIVED] Unified consequence}:
        \begin{center}
            Discrete physical spectra arise dynamically from delay-induced mode selection.
        \end{center}
        
        This replaces operator-based quantization with a dynamical mechanism.
    
    \subsection{Charge and Circulation}
        
        \textbf{[SPECULATIVE] Provisional bridge only.} The canon retains circulation as a structurally important invariant, but it does \emph{not} yet claim a completed derivation of observed gauge charges from circulation data alone. At most, the present document records a working hypothesis that orientation and magnitude of circulation may enter a later charge-assignment program.
    
    \subsection{Unification of Interactions}
        
        The canon records a \emph{unification program}, not yet a completed unification theorem. In the present document:
        \begin{itemize}
            \item electromagnetic effects are represented by the minimal transverse-field bridge of Section~\ref{sec:embridge},
            \item gravitational effects are represented by the clock-field and Poisson-limit bridge of Section~\ref{sec:time},
            \item inertial mass is organized by the medium/topology/geometry/clock structure of Eq.~\eqref{eq:sst_master_equation}.
        \end{itemize}
        
        \textbf{[SPECULATIVE] Programmatic claim:} these sectors may ultimately be read as different effective descriptions of one substrate, but their full common derivation remains open.
    
    \subsection{Atomic Structure as an Emergent Layer}
        
        From Section~\ref{sec:atomic}:
        
        \begin{itemize}
            \item Standard quantum mechanics appears as an effective description
            \item SST introduces circulation-dependent corrections
        \end{itemize}
        
        \textbf{[SPECULATIVE] Unified interpretation}:
        \begin{center}
            Quantum mechanics is an emergent effective theory of underlying vortex dynamics.
        \end{center}
    
    \subsection{Consistency with Spectroscopy}
        
        From Section~\ref{sec:spectroscopy}:
        
        \begin{itemize}
            \item Internal modes are suppressed by an excitation gap
            \item Low-energy physics remains unchanged
        \end{itemize}
        
        \textbf{[ORTHODOX] Constraint}:
        The unified picture does not violate known experimental constraints.
    
    \subsection{Hierarchy of Description}
        
        The theory admits multiple levels:
        
        \begin{enumerate}
            \item \textbf{Microscopic level:} vortex filament dynamics
            \item \textbf{Mesoscopic level:} delay-induced mode structure
            \item \textbf{Macroscopic level:} effective quantum description
        \end{enumerate}
    
    \subsection{Conceptual Summary}
        
        The SST framework can be summarized as:
        
        \begin{center}
            \begin{tabular}{ll}
                \textbf{Mass} & $\rightarrow$ trapped circulation energy \\
                \textbf{Time} & $\rightarrow$ relational swirl-clock field \\
                \textbf{Charge} & $\rightarrow$ provisional circulation bridge \\
                \textbf{Quantization} & $\rightarrow$ delay-induced stability \\
            \end{tabular}
        \end{center}
    
    \subsection{Canonical Statement}
        
        \begin{center}
            \textit{
                All observed physical structure emerges from the dynamics of a continuous medium in which stable vortex configurations, governed by quantized circulation and delay dynamics, produce the phenomena conventionally attributed to particles, fields, and quantum mechanics.
            }
        \end{center}
    
    \subsection{Scope and Limits}
        
        \textbf{[IMPORTANT]}:
        
        \begin{itemize}
            \item The framework is not yet a complete fundamental theory
            \item Several mechanisms remain to be derived rigorously
            \item Many results are conditional or phenomenological
        \end{itemize}
    
    \subsection{Final Consistency Principle}
        
        \begin{center}
            \textit{
                A valid unified interpretation must preserve all experimentally verified results while providing a deeper structural explanation of their origin.
            }
        \end{center}

    \section{Integration of Prior Canon Versions}
        \label{sec:integration}

    \subsection{Integration policy (v0.5.12 reintegration)}

        This block records only material from earlier SST canon layers that can be retained inside the v0.8.1 architecture without weakening the explicit classification discipline of Section~\ref{sec:consistency}.

        Older canon documents are therefore \emph{not} cited as authorities. Transferable content is rewritten here as current-canon statements, each labeled according to its present status.

    \subsection{Retained v0.5.12 kernel sectors}

        The present draft now retains the following v0.5.12 kernel sectors in rewritten form:
        \begin{itemize}
            \item the Chronos--Kelvin invariant and clock--radius transport law (Section~\ref{subsec:chronos_kelvin});
            \item the reparametrized zero-parameter corollary and coarse-graining rule (Section~\ref{subsec:zero_parameter_corollary});
            \item a minimal Swirl--EM bridge and photon / unknot sector (Section~\ref{sec:embridge});
            \item the swirl pressure law as the Euler-level force bridge (Section~\ref{subsec:swirl_pressure_law});
            \item the analogue-metric / clock-field bridge (Section~\ref{sec:time}).
        \end{itemize}

    \subsection{Wave--particle phase rule and measurement bridge}
        \label{subsec:measurement_bridge}

        A retained qualitative rule from earlier canon layers is that SST distinguishes two limiting phases:
        \begin{align}
            \text{R-phase} &:\quad \text{extended, unknotted, radiative}, \\
            \text{T-phase} &:\quad \text{localized, knotted, tangible}.
        \end{align}

        A minimal dynamical bookkeeping law for phase conversion is

        \begin{align}
            \partial_t P_T
            =
            \Gamma_{R\to T}[\mathcal{K}]\,P_R
            -
            \Gamma_{T\to R}[\mathcal{K}]\,P_T,
            \qquad
            P_R + P_T = 1,
            \label{eq:rt_phase_rate_equation}
        \end{align}

        where $\mathcal{K}$ denotes an environment- and coherence-dependent kernel functional. The intended interpretation is that ``measurement'' corresponds not to an axiomatically imposed collapse, but to a medium-induced transfer of support from the extended R-phase to the localized T-phase.

        \textbf{[ORTHODOX LIMIT]}: in weak coupling and after coarse graining, Eq.~\eqref{eq:rt_phase_rate_equation} should reduce to ordinary environment-induced decoherence models \cite{Zurek2003}.

        \textbf{[SPECULATIVE]}: the identification of decoherence with an actual R$\leftrightarrow$T phase transition is SST-specific.

    \subsection{Gauge and weak-mixing bridge}
        \label{subsec:gauge_weakmix_bridge}

        The retained gauge-level statement is that the effective low-energy algebra should reduce to

        \begin{align}
            \mathfrak{g}_{\mathrm{eff}} = \mathfrak{su}(3)\oplus\mathfrak{su}(2)\oplus\mathfrak{u}(1),
            \label{eq:effective_gauge_algebra}
        \end{align}

        with gauge potentials interpreted as holonomy data of coarse-grained multi-director transport.

        A minimal bridge ansatz for the weak mixing angle is

        \begin{align}
            \sin^2\theta_W^{\mathrm{eff}} = \frac{K_Y}{K_Y + K_W},
            \label{eq:weak_mixing_bridge}
        \end{align}

        where $K_Y$ and $K_W$ are effective Abelian and weak-sector stiffness scales. Eq.~\eqref{eq:weak_mixing_bridge} is not a completed derivation, but it preserves the v0.5.12 idea that weak mixing should descend from a ratio of medium-response coefficients rather than appear as an unstructured free number.

        \textbf{[ORTHODOX]}: $\theta_W$ is the standard electroweak mixing parameter of the SM \cite{Weinberg1967,PDG2024}.

        \textbf{[SPECULATIVE]}: Eq.~\eqref{eq:weak_mixing_bridge} is an SST bridge ansatz only.

    \subsection{What the v0.5.12 reintegration does not import}

        The present reintegration does \emph{not} import v0.5.12 cosmogony, engineering subsystems, cosmological-term closures, or highly application-specific protocols into the main canon body. Those remain suitable for appendices or standalone papers until they are more tightly tied to the minimal axiom system.

    \subsection{Purpose of this integration layer}
        Version~0.8.1 is not a fresh theory written from zero, but an explicit consolidation of earlier canon lines. This section records what has been retained, what has been renamed or reorganized, and what has been demoted from canon status to research-track status.

    \subsection{Retained from the early canon layers}
        From the v0.1.x--v0.3.x line, the present canon retains the primitive closure logic:
        \begin{itemize}
            \item a continuum substrate with effective density $\rhoF$,
            \item a core scale $r_c$ fixed by Compton anchoring,
            \item circulation quantization through $\Gamma_0 = 2\pi \rc \, \vnorm$,
            \item the interpretation of matter as stable organized structures rather than point primitives.
        \end{itemize}
        These elements remain non-negotiable because later sectors depend on them directly.

    \subsection{Retained from the middle canon layers}
        From the v0.4.x--v0.6.x line, the canon retains three durable structures:
        \begin{enumerate}
            \item the Rosetta strategy for mapping SST statements into orthodox language,
            \item the photon/torsion sector as the minimal radiation bridge,
            \item delay-induced mode selection as the main route to discreteness without operator postulates.
        \end{enumerate}
        These are now distributed across Sections~\ref{sec:delay}, \ref{sec:time}, and \ref{sec:unification} rather than treated as isolated notes.

    \subsection{Retained from the late canon layers}
        From the v0.7.x line, v0.8.1 retains the strongest structurally reusable sectors:
        \begin{itemize}
            \item the Swirl-Clock as the canonical relational-time field,
            \item the atomic bridge as a benchmark consistency layer rather than a complete atomic derivation,
            \item the spectroscopic gap logic as a hard consistency filter,
            \item the topology-first mass program, with topology separated from thermodynamic or radiative dressing,
            \item the EFT-style foliation and gauge-emergence program as a long-range development track.
        \end{itemize}

    \subsection{Notation and density cleanup}
        Earlier canon versions sometimes conflated background density, mass-equivalent density, and saturated core density. The present version adopts the cleaned hierarchy
        \begin{align}
            \rhoF \quad \text{(effective fluid density)},
            \qquad
            \rhoE \quad \text{(swirl energy density)},
            \qquad
            \rhoM = \rhoE/c^2,
            \qquad
            \rhocore \quad \text{(saturated core density)}.
        \end{align}
        Any older passage that identifies a single symbol with more than one of these roles is superseded by the present separation.

    \subsection{Topology-first reinterpretation of layers}
        A major consolidation from the recent lepton and knot-mass program is the separation between topology and layering. Knot type is retained as the primary identity label, while Golden-layer or thermal ladder indices are treated as dressing or coherence variables rather than species-defining labels. This prevents the canon from over-reading numerical layer fits as ontological identities.

    \subsection{Status of optional sectors}
        The following sectors are retained as important but non-primitive research layers:
        \begin{itemize}
            \item full gauge emergence and charge assignments,
            \item dual-vacuum phenomenology and Unruh-echo interpretation,
            \item complete many-electron atomic closure,
            \item full weak-sector and collider phenomenology,
            \item explicit topology-to-particle dictionary beyond the best-supported benchmark assignments.
        \end{itemize}
        They remain in the canon as organized programs, not as closed theorem-level results.

    \subsection{Integration principle}
        The purpose of v0.8.1 is therefore not to maximize novelty per page, but to minimize duplication and to expose the dependency graph of the theory. Whenever two earlier canon fragments make the same structural claim in different language, the present version keeps one canonical form and demotes the others to historical context.

    \subsection{Research-track companion}
        The sectors that are useful for continuity but too long or too speculative for a minimal main-canon reading are also collected in the companion file \texttt{canon-0.8.1-research-track.tex}. That excerpt includes the thermodynamic-radiation interface, the minimal-coupling QED analogy, the ribbon-invariant chirality refinement, the particle-candidate mapping layer, the hydrodynamic exchange / Pauli-barrier estimate, the numerical-benchmark policy block, the gauge-sector roadmap, and the \emph{dark-sector and galactic canonization execution package} (frozen topology toolkit, benchmark protocol template, Euler-pressure anchor details, galactic branch gate, and probe--equation map). The integration subsections below retain the full in-context treatments with labels and cross-references; the companion file provides a portable standalone copy. Editorial separation does not reject these sectors; it prevents the core narrative from over-claiming closure.

    \subsection{Dark-sector and galactic canonization status (v0.8.x)}
        \label{sec:dark-canonicalization-status}

        This block records promotion labels for the dark-sector / galactic program consistent with the internal canonization plan. Detailed freezes, the mandatory benchmark-protocol checklist, the galactic profile branch gate (Branch~A: legacy-profile refine with benchmark traceability), and the probe falsification table live in \texttt{canon-0.8.1-research-track.tex} (Section titled \texttt{Dark-sector and galactic canonization execution package}).

        \begin{table}[h]
        \centering
        \caption{Master tracker: dark-sector and galactic program (v0.8.x).}
        \label{tab:dark-sector-master-tracker}
        \begin{tabular}{p{4.2cm}p{2.3cm}p{2.5cm}p{5.5cm}}
        \hline
        Item & Prior state & Target & Status after Stage 0--4 freeze \\
        \hline
        Dark-sector Euler pressure anchor & BRIDGE & CANON & CANON-CANDIDATE (derivation, assumptions, limits, falsifiers frozen; see Section~\ref{subsec:swirl_pressure_law}) \\
        Dark taxonomy & HOLD & CANON & BRIDGE (symmetry-first classifier + filters; Section~\ref{subsec:dark_knot_classification}) \\
        Galactic swirl law & BRIDGE & CANON & BRIDGE (Branch~A closure; companion execution package) \\
        Experimental probes & HOLD & CANON-BRIDGE & CANON-BRIDGE (equation / null-result map; companion) \\
        Benchmark / reproducibility & BRIDGE & CANON & CANON-CANDIDATE (protocol template frozen; repository-wide reuse pending) \\
        Topological reference identities & CANON-CANDIDATE & CANON & CANON-ready scope/notation (Appendix~\ref{sec:appendix_topology_toolkit}) \\
        \hline
        \end{tabular}
        \end{table}

        \textbf{[DERIVED] Stage 5 review:} cross-item dependencies are respected (pressure anchor and benchmark layer precede full galactic closure; taxonomy strengthens other sectors). No sector is promoted on fit quality alone; each carries explicit falsifiability or reproducibility hooks as in Table~\ref{tab:dark-sector-master-tracker} and the companion execution package.

    \subsection{Integration Policy}

        This section records only material from prior internal SST development that can be retained inside the
        v0.8.1 architecture \emph{without} weakening the epistemic classification framework of
        Section~\ref{sec:consistency}.

        \textbf{Rule.}
        Older canon documents are \emph{not} cited as external authorities. Instead, transferable content is
        rewritten here as current-canon statements, each labeled as \textbf{[ORTHODOX]}, \textbf{[DERIVED]},
        or \textbf{[SPECULATIVE]}.

    \subsection{Selective Recovery from Pre-v0.7 Canon Layers}

        The present draft absorbs only those legacy ingredients that can be rewritten in the current epistemic format without appealing to earlier canon versions as authorities.

        \paragraph{Analogue line-element bridge.}
        Legacy SST drafts recorded a schematic analogue line element and a metric--clock identification of the form
        \begin{align}
            ds^2 &= -c^2\,\SwirlClock(x)^2\,dt^2 + \delta_{ij}\,dx^i\,dx^j,
            \label{eq:legacy_metric}\\
            g_{00}(x) &= -\SwirlClock(x)^2 .
            \label{eq:metric_swirlclock_bridge}
        \end{align}

        Recovered legacy sectors include:
        \begin{itemize}
            \item the finite-core atomic bridge, Eqs.~\eqref{eq:lambda_sst}--\eqref{eq:softcore_potential};
            \item the analogue line-element bridge, Eqs.~\eqref{eq:legacy_metric}--\eqref{eq:metric_swirlclock_bridge};
            \item compact research-track interfaces to thermodynamic radiation and minimal-coupling QED analogies;
            \item topological chirality refinements based on standard ribbon invariants.
        \end{itemize}

        Legacy material not imported includes version history, cosmogonic narrative, and phenomenological galaxy-scale laws that are not yet sufficiently integrated with the present minimal axiom system.

    \subsection{Thermodynamic Radiation Interface}

        A recurring legacy idea is that a local swirl-energy density may define an effective temperature field and thereby an emergent radiation sector.

        Let
        \begin{align}
            T_{\mathrm{eff}} = T_{\mathrm{eff}}\!\left(\rho_E,\,\omega,\,S_{(t)}\right)
            \label{eq:teff_placeholder}
        \end{align}
        denote an unspecified constitutive relation. A minimal legacy interface then asks whether Planck-type spectra can emerge from a medium whose local excitation scale is set by $T_{\mathrm{eff}}$.

        \textbf{[ORTHODOX]}: blackbody spectra are governed by Planck's law once a valid temperature variable is defined.\newline
        \textbf{[SPECULATIVE]}: SST does not yet derive Eq.~\eqref{eq:teff_placeholder}, so this sector remains a research program rather than canonically closed physics.

    \subsection{Minimal Coupling Interface to Quantum Electrodynamics}

        Early research notes repeatedly explored whether an effective vector potential could be built from circulation variables. The most conservative placeholder takes the form
        \begin{align}
            \nabla \rightarrow \nabla - i\,\frac{m}{\hbar}\,\mathbf{A}_{\mathrm{eff}},
            \qquad
            \mathbf{A}_{\mathrm{eff}} = \chi\,\mathbf{v},
            \label{eq:legacy_minimal_coupling}
        \end{align}
        where $\chi$ is a scalar weight and $\mathbf{v}$ is the local carrier velocity field.

        \textbf{[ORTHODOX]}: minimal coupling is the standard gateway from free to gauge-coupled wave dynamics.\newline
        \textbf{[SPECULATIVE]}: Eq.~\eqref{eq:legacy_minimal_coupling} is only an analogy template unless a full gauge-covariant derivation is supplied.

    \subsection{Ribbon-Invariant Chirality Refinement}

        The older canon layers also suggested using the C\u{a}lug\u{a}reanu--White--Fuller relation
        \begin{align}
            Lk = Tw + Wr
            \label{eq:calugareanu}
        \end{align}
        as a refinement of the knot-based matter/antimatter dictionary.

        \textbf{[ORTHODOX]}: Eq.~\eqref{eq:calugareanu} is a standard geometric identity for framed curves and ribbons.\newline
        \textbf{[DERIVED]}: it can be used internally to separate writhe-dominated from twist-dominated realizations of the same coarse knot class.\newline
        \textbf{[SPECULATIVE]}: the mapping from sign$(Tw)$ or related chirality measures to particle/antiparticle sectors remains a model hypothesis.

    \subsection{Particle-Candidate Mapping Layer}
        \label{subsec:integration_knotmap}

        \textbf{[SPECULATIVE]} Working dictionary: a minimal knot / composition mapping retained from earlier SST work is summarized
        in Table~\ref{tab:particle_candidate_mapping}. This table is organizational rather than definitive; it records
        the current topological assignments used by the mass and taxonomy program.

        \begin{table}[h]
        \centering
        \caption{Minimal knot-class to particle-candidate summary retained in v0.8.1.}
        \label{tab:particle_candidate_mapping}
        \begin{tabular}{lll}
        \hline
        Knot class / composition & Candidate & Status / role \\
        \hline
        $3_1$ & electron & baseline lepton calibration \\
        $T(2,5)$ or equivalent higher torus class & muon & lepton hierarchy candidate \\
        $T(2,7)$ or equivalent higher torus class & tau & lepton hierarchy candidate \\
        $5_2 + 5_2 + 6_1$ & proton & composite baryon candidate \\
        analogous nearby composite class & neutron & composite baryon candidate \\
        \hline
        \end{tabular}
        \end{table}

        \textbf{[CRITICAL NOTE]} Table~\ref{tab:particle_candidate_mapping} is a working dictionary only. A full canonical
        assignment requires an explicit invariant set and a mass-functional comparison program.

    \subsection{Hydrodynamic Exchange and the Pauli Barrier}
        \label{subsec:integration_pauli}

        \textbf{[ORTHODOX]} Fluid-mechanical input: for two thin filamentary loops \(C_1\) and \(C_2\) in an
        incompressible medium, the Biot--Savart interaction energy is \cite{Batchelor1967}
        \begin{align}
            E_{\mathrm{int}}
            =
            \frac{\rho_{\!f}\Gamma_1\Gamma_2}{4\pi}
            \oint_{C_1}\oint_{C_2}
            \frac{d\mathbf{l}_1\cdot d\mathbf{l}_2}
            {\lVert \mathbf{r}_1 - \mathbf{r}_2 \rVert}.
            \label{eq:biot_savart_interaction}
        \end{align}

        \textbf{[DERIVED]} SST bridge: if two identical electron-candidate loops are forced into the same spatial state,
        regularization at the core radius \(r_c\) yields the overlap penalty
        \begin{align}
            V_{\mathrm{Pauli}}(\mathbf{d})
            \approx
            \frac{\rho_{\!f}\Gamma_0^2}{4\pi}
            \oint\!\!\oint
            \frac{d\mathbf{l}_1\cdot d\mathbf{l}_2}
            {\sqrt{\lVert \mathbf{r}_1(\theta_1)-\mathbf{r}_2(\theta_2)+\mathbf{d}\rVert^2+r_c^2}},
            \label{eq:pauli_barrier_regularized}
        \end{align}
        with maximal-overlap scale
        \begin{align}
            V_{\mathrm{Pauli}}^{\max}
            \approx
            \frac{\rho_{\!f}\Gamma_0^2}{4\pi}\left(\frac{L}{r_c}\right)\mathcal{O}(1).
            \label{eq:pauli_barrier_scale}
        \end{align}

        Using the canonical values
        \begin{align}
            \rho_{\!f} &= 7.0 \times 10^{-7}\ \mathrm{kg\,m^{-3}}, \nonumber\\
            r_c &= 1.40897017 \times 10^{-15}\ \mathrm{m}, \nonumber\\
            \Gamma_0 &= 2\pi r_c \lVert \mathbf{v}_{\!\boldsymbol{\circlearrowleft}} \rVert
            \approx 9.6836\times 10^{-9}\ \mathrm{m^2\,s^{-1}},
        \end{align}
        and \(L \sim 2\pi a_0\) at the hydrogenic scale, one obtains the benchmark estimate
        \begin{align}
            V_{\mathrm{Pauli}}^{\max}
            \sim 1.23\times 10^{-18}\ \mathrm{J}
            \sim 7.6\ \mathrm{eV}.
            \label{eq:pauli_barrier_numeric}
        \end{align}

        \textbf{[CRITICAL NOTE]} This is retained as a canonical scale benchmark, not yet as a closed theorem. Its role is
        to show that the exclusion-energy estimate falls in an atomically relevant window rather than at obviously
        pathological scales.

        As written, Eq.~\eqref{eq:pauli_barrier_numeric} should be interpreted as an order-of-magnitude benchmark attached to the canonical constants, not yet as a unique parameter-free prediction.

    \subsection{Numerical Benchmark Layer}
        \label{subsec:integration_benchmarks}

        \textbf{[DERIVED]} Computational-program requirement: earlier SST work contained explicit benchmark tables for particle and
        atomic masses. In v0.8.1, the benchmark layer is retained as a reproducibility requirement rather than as a
        standalone authority claim.

        \begin{table}[h]
        \centering
        \caption{Benchmark summary structure to be populated from canonical scripts and archived CSV output.}
        \label{tab:benchmark_summary_structure}
        \begin{tabular}{lrrrrc}
        \hline
        Particle & $m_{\mathrm{exp}}$ (MeV) & $m_{\mathrm{SST}}$ (MeV) & rel.\ error & mode & class \\
        \hline
        $e$   & 0.51099895 & 0.51099895 & 0 & predictive/closure & $3_1$ \\
        $\mu$ & 105.6583745 & \dots & \dots & predictive & candidate torus class \\
        $\tau$& 1776.86 & \dots & \dots & predictive & candidate torus class \\
        $p$   & 938.272088 & \dots & \dots & predictive or closure & composite class \\
        $n$   & 939.565420 & \dots & \dots & predictive or closure & composite class \\
        \hline
        \end{tabular}
        \end{table}

        \textbf{[CRITICAL NOTE]} A future v0.8.x benchmark appendix should include:
        \begin{itemize}
            \item the exact script or package path used for generation,
            \item archived CSV outputs,
            \item mode definitions (predictive vs.\ exact-closure),
            \item uncertainty propagation on canonical constants and geometric factors,
            \item theorem-style pass/fail metadata (including tail-quality checks, postchecks, and extrapolation summaries),
            \item explicit separation between practical best-estimate closure and theorem-level closure,
            \item normalization of legacy constant names into SST notation (\(\rhoF,\rc,\Gamma_0\)) while preserving numerical values.
        \end{itemize}

        This requirement is included to keep the benchmark layer falsifiable and reproducible, rather than merely illustrative.

    \subsection{Gauge-Sector Roadmap}
        \label{subsec:integration_gauge}

        \textbf{[SPECULATIVE]} Structural roadmap: a retained minimal statement of the gauge program is:
        \begin{align}
            \mathfrak{g}_{\mathrm{eff}}
            =
            \mathfrak{su}(3)\oplus\mathfrak{su}(2)\oplus\mathfrak{u}(1),
        \end{align}
        with effective gauge phases organized as holonomy data of multi-director transport.

        A minimal phase-based representation is
        \begin{align}
            A_\mu^{\,a}(x) = \partial_\mu \theta^a(x),
            \label{eq:gauge_phase_potential}
        \end{align}
        and clock coupling acts as a common phase reference,
        \begin{align}
            \theta^a(x) \mapsto \theta^a(x) + q^a \chi(x)
            \quad\Rightarrow\quad
            A_\mu^{\,a} \mapsto A_\mu^{\,a} + q^a \partial_\mu \chi.
            \label{eq:gauge_clock_shift}
        \end{align}

        \textbf{[CRITICAL NOTE]} The gauge layer is retained in v0.8.1 only as a roadmap statement. Running couplings,
        anomaly cancellation, and realistic mixing matrices remain open work.

    \section{Discussion and Limits}
        \label{sec:discussion}

    \subsection{Status of the recovered v0.5.12 sectors}

        The material reintroduced here does not have uniform epistemic status.

        \begin{itemize}
            \item The Chronos--Kelvin and pressure-law sectors are orthodox fluid-mechanical structures with SST reinterpretations.
            \item The analogue-metric and clock-field EFT layers are orthodox comparison tools, but only partially derived within SST.
            \item The R/T measurement bridge and weak-mixing stiffness ansatz remain explicitly speculative.
        \end{itemize}

    \subsection{What v0.8.1 now recovers from v0.5.12}

        After the present patch, v0.8.1 no longer omits the main v0.5.12 kernel sectors that are structurally compatible with the current architecture:
        \begin{itemize}
            \item Chronos--Kelvin / clock transport,
            \item zero-parameter reparameterization discipline,
            \item Swirl--EM and photon / unknot bridge,
            \item swirl pressure law,
            \item analogue metric and clock-field gravity bridge,
            \item R/T measurement bridge,
            \item weak-mixing bridge ansatz.
        \end{itemize}

    \subsection{Open items that remain outside the present closure}

        The following remain open and are \emph{not} upgraded to full canonical closure in this draft:
        \begin{itemize}
            \item a first-principles derivation of the full SM gauge sector;
            \item a completed EWSB derivation tied to the canon constants;
            \item a rigorous derivation of the R/T measurement kernel from continuum dynamics;
            \item full CANON promotion of the dark-sector taxonomy and galactic closure beyond the current BRIDGE labels (Section~\ref{sec:dark-canonicalization-status}).
        \end{itemize}

    \subsection{What is established at canon level}
        The present document supports the following as canon-level structures:
        \begin{itemize}
            \item the primitive constant chain $(\rhoF, v_{\swirlarrow}, \omega_c) \to r_c \to \Gamma_0$,
            \item the Swirl-Clock as the coarse-grained relational time factor,
            \item delay-induced mode selection as the preferred discreteness mechanism,
            \item the atomic bridge and spectroscopic sectors as compatibility filters,
            \item the topology-first program for mass organization.
        \end{itemize}
        These claims are not equally proven, but they are the sectors on which the rest of the canon is structurally built.

    \subsection{What remains closure-level or conjectural}
        Several central SST ambitions remain open:
        \begin{enumerate}
            \item a rigorous derivation of full particle spectra from topology alone,
            \item a first-principles derivation of charge assignments and gauge couplings,
            \item a complete many-body derivation of shell filling and exclusion,
            \item a full demonstration that delay plus topology reproduces the observed quantum rule set without importing it.
        \end{enumerate}
        The canon records these as active programs rather than as settled conclusions.

    \subsection{Compatibility limits}
        Any viable SST realization must remain compatible with:
        \begin{itemize}
            \item atomic spectroscopy,
            \item relativistic effective kinematics,
            \item precision electroweak data where an SST bridge is claimed,
            \item the absence of large low-energy deviations from standard quantum mechanics.
        \end{itemize}
        This is why the canon repeatedly imposes suppression, weak-coupling, and low-energy recovery conditions.

    \subsection{Minimal falsifiers}
        The present canon would be materially weakened by any of the following:
        \begin{itemize}
            \item failure of the constant chain to remain dimensionally and numerically self-consistent,
            \item absence of any workable suppression mechanism for internal modes at atomic energies,
            \item inability to formulate a non-circular route from delay dynamics to discrete stable spectra,
            \item systematic contradiction between any SST atomic bridge and precision spectroscopy.
        \end{itemize}
        Conversely, strong support would come from a successful theorem ladder showing how discrete spectra, topological stabilization, and low-energy effective quantum structure arise from one shared substrate model.

    \subsection{Methodological limit}
        The canon is deliberately hybrid in style. Some sectors are written as exact closures, some as effective models, and some as programmatic roadmaps. This is a feature rather than a flaw, provided the epistemic status of each claim remains explicit. The document should therefore not be judged as though every section were intended to be a stand-alone journal theorem.

    \subsection{Near-term priorities}
        The next canon-strengthening tasks are clear:
        \begin{enumerate}
            \item complete the theorem ladder for the mass and clock sectors,
            \item expand numerical validation blocks using the canonical constants,
            \item refine the atomic bridge into sharper benchmark observables,
            \item decide which optional sectors belong in the core canon and which should migrate to separate technical notes or the research-track companion.
        \end{enumerate}

    \subsection{Final limit statement}
        SST is presently best understood as a structured unification program with several strong internal bridges and several open derivational gaps. The canon is therefore mature enough to organize the theory coherently, but not yet so mature that every phenomenological sector should be presented as theorem-complete.

    \subsection{Status of Legacy Recoveries}

        The additional material imported here from older canon layers does not have uniform status.

        \begin{itemize}
            \item The analogue-metric bridge and ribbon identity are mathematically orthodox tools.
            \item Their SST interpretation is partially derived but not fully closed dynamically.
            \item The thermal-radiation and QED-interface sectors remain explicitly research-track.
        \end{itemize}

        This section therefore serves as a controlled retention layer, not as a blanket validation of all legacy material.

        Accordingly, these recoveries should be read as structured extensions of the minimal core, not as fully promoted final doctrine.

    \subsection{What v0.8.1 Now Retains}

        After integration, the present Canon retains the following prior-development sectors in current-canon form:
        \begin{itemize}
            \item a particle-candidate topological dictionary (Section~\ref{subsec:integration_knotmap});
            \item a hydrodynamic exchange-energy estimate for fermionic exclusion
            (Section~\ref{subsec:integration_pauli});
            \item a benchmark / reproducibility layer (Section~\ref{subsec:integration_benchmarks});
            \item a minimal gauge-sector roadmap (Section~\ref{subsec:integration_gauge});
            \item a clock-field EFT bridge and Poisson gravity limit (Section~\ref{sec:time}).
        \end{itemize}

    \subsection{What is Deliberately Not Claimed}

        The present document still does \emph{not} claim:
        \begin{itemize}
            \item a completed first-principles derivation of the Standard Model;
            \item a unique finalized knot-to-particle dictionary;
            \item a closed proof of the exclusion barrier beyond the current scaling benchmark;
            \item a fully renormalized relativistic field theory for the gauge and clock sectors.
        \end{itemize}

        These omissions are deliberate. v0.8.1 is intended to remain logically strict even when phenomenological
        sectors are reintroduced.

    \subsection{Canon Rule Going Forward}

        Any further reintegrated material must satisfy all of the following:
        \begin{enumerate}
            \item it can be written without citing older internal canon versions as authorities;
            \item it can be assigned an explicit epistemic label;
            \item it includes either an internal derivation path or a reproducible benchmark route;
            \item it preserves compatibility with the orthodox limits recorded earlier in the document.
        \end{enumerate}

% ================================================================================= %
% APPENDICES
% ================================================================================= %
\newpage
\section*{Appendices}
    \appendix

    \section{Topological reference toolkit}
        \label{sec:appendix_topology_toolkit}

        The following orthodox items form the mandatory internal mathematical toolkit cited by topology-first SST sectors. Statements are standard; SST-specific mappings appear in the main text with explicit status labels.

        \begin{enumerate}
            \item Gauss linking integral and linking number for disjoint closed curves.
            \item Helicity $\mathcal{H}=\int \mathbf{v}\cdot(\nabla\times\mathbf{v})\,dV$ in ideal (inviscid) fluid flow and its conservation in the appropriate limit.
            \item Thin-tube reduction of helicity to flux--linkage form in the standard filament limit.
            \item Hopf invariant for localized field configurations (degree--type integral) where used as a topological density.
            \item Basic torus-knot invariants as needed in the mass program (e.g.\ minimal crossing number, braid index, genus) with fixed naming consistent with Rolfsen notation in the main text.
            \item C\u{a}lug\u{a}reanu--White--Fuller relation $Lk=Tw+Wr$ for framed ribbons (already recalled as Eq.~\eqref{eq:calugareanu} in the integration layer).
        \end{enumerate}

        \textbf{[ORTHODOX]}: each item is classical mathematics \cite{Moffatt1969,White1969,Batchelor1967}. \textbf{[DERIVED]}: scope and notation for v0.8.x are frozen to this list plus house symbols $(\rhoF,\Gamma,\rc,\mathcal{H})$ without parallel aliases.

    \section{Full Derivations}
        \label{sec:appendix_derivations}

        \subsection{Canonical constant chain}
            The present canon uses the closure chain
            \begin{align}
                r_e &= \frac{\alpha \hbar}{m_e c}, \\
                r_c &= \frac{r_e}{2}, \\
                \omega_c &= \frac{m_e c^2}{\hbar}, \\
                v_{\swirlarrow} &= \omega_c r_c, \\
                \Gamma_0 &= 2\pi r_c v_{\swirlarrow}.
            \end{align}
            Substituting $r_c = \alpha\hbar/(2m_e c)$ into $v_{\swirlarrow}=\omega_c r_c$ gives
            \begin{align}
                v_{\swirlarrow}
                =
                \left(\frac{m_e c^2}{\hbar}\right)
                \left(\frac{\alpha \hbar}{2m_e c}\right)
                = \frac{\alpha}{2}c,
            \end{align}
            hence the canonical ratio
            \begin{align}
                \eta_0 = \frac{v_{\swirlarrow}}{c} = \frac{\alpha}{2}.
            \end{align}

        \subsection{Core density closure}
            Writing the electron rest energy as localized swirl energy inside a core region leads to the closure
            \begin{align}
                \rho_{\mathrm{core}}
                \sim
                \frac{m_e c^2}{2\pi v_{\swirlarrow}^2 r_c^3},
            \end{align}
            which is the form already recorded in Eq.~\eqref{eq:core_density}. This relation is dimensionally exact:
            \begin{align}
                \left[\frac{m_e c^2}{v^2 r^3}\right]
                =
                \frac{\mathrm{kg\,m^2\,s^{-2}}}{\mathrm{m^2\,s^{-2}}\,\mathrm{m^3}}
                = \mathrm{kg\,m^{-3}}.
            \end{align}

        \subsection{Geometric gate identity}
            Using the Compton wavelength $\lambda_c = h/(m_e c)$ and the core closure $r_c = \alpha\hbar/(2m_e c)$, one obtains
            \begin{align}
                \frac{\lambda_c}{\pi r_c}
                =
                \frac{h/(m_e c)}{\pi \alpha \hbar /(2m_e c)}
                =
                \frac{2h}{\pi \alpha \hbar}
                =
                \frac{4}{\alpha}.
            \end{align}
            This is the canonical geometric shielding factor used throughout the late mass-functional papers.

        \subsection{Useful numerical anchors}
            With the canonical values
            \begin{align}
                v_{\swirlarrow} &\approx 1.09384563\times 10^6\ \mathrm{m\,s^{-1}}, \\
                r_c &\approx 1.40897017\times 10^{-15}\ \mathrm{m}, \\
                \Gamma_0 &\approx 9.6836\times 10^{-9}\ \mathrm{m^2\,s^{-1}},
            \end{align}
            one has the Compton turnover time
            \begin{align}
                \tau_c = \frac{2\pi \rc}{\vnorm} = \frac{2\pi}{\omega_c},
            \end{align}
            and the geometric gate
            \begin{align}
                \frac{\lambda_c}{\pi \rc} \approx 5.48\times 10^2.
            \end{align}

    \section{Delay Mathematics}
        \label{sec:appendix_delay}

        \subsection{Fixed-point condition for locked modes}
            Starting from the delayed phase equation
            \begin{align}
                \dot{\phi}(t) = \omega_0 + \kappa \sin\big(\phi(t-\tau)-\phi(t)\big),
            \end{align}
            seek a rigidly rotating solution
            \begin{align}
                \phi(t)=\Omega t + \phi_\star.
            \end{align}
            Substitution gives the algebraic mode condition
            \begin{align}
                \Omega = \omega_0 + \kappa \sin(\Omega\tau),
            \end{align}
            which admits multiple branches labelled by the effective phase winding across the delay interval.

        \subsection{Linear stability of a locked branch}
            Let
            \begin{align}
                \phi(t) = \Omega t + \epsilon(t),
                \qquad |\epsilon|\ll 1.
            \end{align}
            Linearizing yields
            \begin{align}
                \dot{\epsilon}(t)
                =
                \kappa \cos(\Omega\tau)
                \big(\epsilon(t-\tau)-\epsilon(t)\big).
            \end{align}
            Using the exponential ansatz $\epsilon(t)\propto e^{\lambda t}$ gives the characteristic equation
            \begin{align}
                \lambda
                =
                \kappa \cos(\Omega\tau)
                \left(e^{-\lambda\tau}-1\right).
            \end{align}
            Stability requires $\Re(\lambda)<0$ for all nontrivial roots. Thus discrete spectra in SST are not arbitrary: only those branches satisfying both the fixed-point equation and the stability inequality survive dynamically.

        \subsection{Small-delay expansion}
            For $|\lambda\tau|\ll 1$,
            \begin{align}
                e^{-\lambda\tau} = 1 - \lambda\tau + \mathcal{O}((\lambda\tau)^2),
            \end{align}
            so the characteristic equation becomes
            \begin{align}
                \lambda \approx -\kappa\tau\cos(\Omega\tau)\,\lambda.
            \end{align}
            Nontrivial damping therefore requires the effective factor $\kappa\tau\cos(\Omega\tau)$ to have the correct sign. This provides a simple branch-selection criterion at weak delay.

        \subsection{Interpretive consequence}
            The important canon point is that discreteness is produced in two stages:
            \begin{align}
                \text{delay} \;\longrightarrow\; \text{multiple algebraic branches} \;\longrightarrow\; \text{stability filter}.
            \end{align}
            Quantization is therefore not inserted by hand; it is the residue of a dynamical selection problem on a delayed closed carrier.

    \section{Conversation-Derived Insights}
        \label{sec:appendix_conversations}

        \subsection{Role of this appendix}
            This appendix records recurrent insights that emerged during the project-level development of the canon and are useful for guiding future work, even when they do not yet qualify as theorem-level results.

        \subsection{Topology first, layers second}
            A central refinement of the recent canon is the separation
            \begin{align}
                \text{particle identity} \;\longleftrightarrow\; \text{topology},
                \qquad
                \text{excitation or dressing state} \;\longleftrightarrow\; \text{layer index}.
            \end{align}
            This prevents Golden-layer fits from being mistaken for a species ontology and keeps the topology program structurally cleaner.

        \subsection{Spectral origin of \texorpdfstring{$\zeta(3)$}{zeta(3)}}
            The project discussions repeatedly converged on the conclusion already encoded in Section~\ref{subsec:spectral_locked_carriers}: the appearance of $\zeta(3)$ is best interpreted as a spectral invariant of a cubic-weighted closed-carrier mode tower, not as a direct invariant of threefold geometry by itself.

        \subsection{Atomic bridge status discipline}
            Another recurring project insight is that the atomic bridge is strongest when treated as a benchmark layer with explicit status tags. Hydrogenic consistency, branch-sensitive couplings, and the spectroscopic gap form a robust bridge; shell filling, exclusion, and full heavy-atom structure remain open programs rather than completed derivations.

        \subsection{Master-equation viewpoint}
            A major organizational gain of the recent conversations is the recognition that many SST sectors can be read as different projections of one shared structure:
            \begin{align}
                \text{medium scale} \times \text{topology} \times \text{geometry} \times \text{clocking}.
            \end{align}
            Even when this is not written as a single canonical equation inside the present draft, it remains a useful guide for reducing duplication across the mass, time, and thermodynamic sectors.

        \subsection{Dual-vacuum and photon sector status}
            The project-level discussions also clarified a canon policy: the photon/torsion bridge and the dual-vacuum phenomenology are valuable and fertile, but they should not silently override the simpler core canon. They are best treated as phenomenological extensions built on top of the primitive constant chain and Swirl-Clock structure.

        \subsection{Open theorem ladder}
            The most important unfinished task identified across the project is not adding new sectors, but tightening the derivational ladder between existing ones. In practice this means:
            \begin{enumerate}
                \item turning heuristic closures into explicit lemmas,
                \item attaching numerical validation to each nontrivial scaling law,
                \item separating benchmark results from full ontological claims.
            \end{enumerate}
            This appendix therefore serves as a reminder that canon quality increases more by sharpening structure than by increasing conceptual breadth.

% ===========================
% SUPPLEMENTAL APPENDIX
% ===========================
 
        \newpage
    
    
    \section{Explicit Topological Mass Closures}
    \label{sec:appendix_explicit_mass}
    
    \subsection{Purpose and status}
        This appendix supplies an explicit closure layer for the mass program without replacing the canonical master equation of Eq.~\eqref{eq:sst_master_equation}. Its role is to make concrete those invariant choices that were historically used in pre-v0.8 layers whenever actual particle-mass calculations were attempted.
        
        \textbf{[DERIVED]} Canon rule: the abstract kernel \(\Xi_K\) of the main text may be specialized to more explicit invariant combinations provided the resulting branch remains dimensionless, benchmarkable, and subordinate to the topology-first hierarchy.
        
        \textbf{[CRITICAL NOTE]} Nothing in this appendix upgrades the explicit kernel to a theorem of first-principles filament dynamics. The present role is closure discipline and reproducible benchmarking.
    
    \subsection{Explicit invariant branch of the canonical kernel}
        A conservative explicit specialization is obtained by resolving the topological kernel into ropelength, braid complexity, and Seifert-genus data:
        \begin{align}
            \Xi_K^{(\mathrm{exp})}
            =
            \left[
                \alpha_{\mathcal{L}}
                \frac{\mathcal{L}_{\mathrm{tot}}(K)}{\mathcal{L}_{\mathrm{tot}}(3_1)}
                +
                \alpha_b \bigl(b(K)-1\bigr)
            +
                \alpha_g\, g(K)
            \right]
            \varphi^{-2k(K)}.
            \label{eq:explicit_kernel}
        \end{align}
        Here:
        \begin{itemize}
            \item \(\mathcal{L}_{\mathrm{tot}}(K)\) is the total ropelength-like invariant used for the chosen embedding or ideal representative,
            \item \(b(K)\) is the braid index,
            \item \(g(K)\) is the Seifert genus,
            \item \(k(K)\) is the canonical layer / dressing index already used in the main text.
        \end{itemize}
        
        \textbf{[ORTHODOX]} The invariants \(\mathcal{L}_{\mathrm{tot}}(K)\), \(b(K)\), and \(g(K)\) are standard topological or geometric knot descriptors.
        
        \textbf{[DERIVED]} Equation~\eqref{eq:explicit_kernel} is dimensionless and therefore admissible as a branch of the canonical kernel.
        
        Substituting Eq.~\eqref{eq:explicit_kernel} into Eq.~\eqref{eq:sst_master_equation} yields the explicit benchmark mass branch
        \begin{align}
            M_K^{(\mathrm{exp})}(x)
            =
            \rhoM(x)\,\rc^3\,\Pi_K\,\Xi_K^{(\mathrm{exp})}\,\SwirlClock(x)^{-2}.
            \label{eq:explicit_mass_branch}
        \end{align}
    
    \subsection{Baryon benchmark closure}
        For composite baryons built from constituent topology classes \(K_u\) and \(K_d\), a conservative additive benchmark branch is
        \begin{align}
            M_B^{(\mathrm{add})}
            =
            \Pi_B\,M_0
            \left(
                n_u\,\Xi_{K_u}^{(\mathrm{exp})}
                +
                n_d\,\Xi_{K_d}^{(\mathrm{exp})}
            \right),
            \label{eq:baryon_additive_branch}
        \end{align}
        where \(n_u\) and \(n_d\) are constituent counts and \(M_0\) is the common medium scale extracted from the main mass functional.
        
        A retained legacy golden-closure benchmark is
        \begin{align}
            M_p^{(\varphi)}
            &=
            \varphi^{-2} 3^{-1/\varphi}\,\bigl(2M_u + M_d\bigr),
            \label{eq:proton_phi_benchmark}\\
            M_n^{(\varphi)}
            &=
            \varphi^{-2} 3^{-1/\varphi}\,\bigl(M_u + 2M_d\bigr).
            \label{eq:neutron_phi_benchmark}
        \end{align}
        
        \textbf{[SPECULATIVE]} Equations~\eqref{eq:proton_phi_benchmark}--\eqref{eq:neutron_phi_benchmark} are retained as baryon-sector closure benchmarks only. They are not promoted to universal mass laws.
    
    \subsection{Canonical use policy}
        The explicit closures above are intended for:
        \begin{enumerate}
            \item theorem-ladder benchmarking,
            \item archived CSV reproduction,
            \item side-by-side comparison of topology branches,
            \item explicit baryon fits.
        \end{enumerate}
        They are \emph{not} intended to replace the abstract organizing form of the main text.
    
    \section{Effective Quantum Bridge}
    \label{sec:appendix_quantum_bridge}
    
    \subsection{Purpose and status}
        The main canon states that quantum mechanics should appear as an effective low-energy description of the medium. The present appendix supplies the minimal bridge formulas needed to make that statement mathematically legible without claiming a completed derivation of the full quantum rule set.
        
        \textbf{[ORTHODOX]} The Madelung decomposition and Schr\"odinger-form continuum limit are standard.
        
        \textbf{[DERIVED]} SST may use these formulas as a coarse-grained field representation once a local amplitude density and phase have been specified.
        
        \textbf{[SPECULATIVE]} The identification of the resulting fields with branch-resolved swirl carriers remains an SST interpretation.
    
    \subsection{Amplitude--phase decomposition}
        Introduce an effective complex mode field
        \begin{align}
            \psi(x,t) = \sqrt{\rho_\psi(x,t)}\,e^{i\theta(x,t)},
            \label{eq:madelung_ansatz_appendix}
        \end{align}
        with \(\rho_\psi \ge 0\) and real phase \(\theta\). The associated effective transport velocity is
        \begin{align}
            \mathbf{u}_\psi
            =
            \frac{\hbar}{m_{\mathrm{eff}}}\nabla\theta.
            \label{eq:effective_phase_velocity}
        \end{align}
        
        Assuming a Schr\"odinger-form bridge equation
        \begin{align}
            i\hbar \partial_t \psi
            =
            \left(
                -\frac{\hbar^2}{2m_{\mathrm{eff}}}\nabla^2 + V_{\mathrm{eff}}
            \right)\psi,
            \label{eq:schrodinger_bridge_appendix}
        \end{align}
        one obtains the equivalent real system
        \begin{align}
            \partial_t \rho_\psi + \nabla\cdot\bigl(\rho_\psi \mathbf{u}_\psi\bigr) &= 0,
            \label{eq:continuity_bridge_appendix}\\
            \partial_t \theta
            +
            \frac{m_{\mathrm{eff}}}{2}\lVert \mathbf{u}_\psi\rVert^2
            + V_{\mathrm{eff}}
            + Q_{\mathrm{eff}} &= 0,
            \label{eq:hamjac_bridge_appendix}
        \end{align}
        with effective quantum potential
        \begin{align}
            Q_{\mathrm{eff}}
            =
            -\frac{\hbar^2}{2m_{\mathrm{eff}}}
            \frac{\nabla^2 \sqrt{\rho_\psi}}{\sqrt{\rho_\psi}}.
            \label{eq:quantum_potential_appendix}
        \end{align}
        
        \textbf{[ORTHODOX]} Equations~\eqref{eq:continuity_bridge_appendix}--\eqref{eq:quantum_potential_appendix} are the standard Madelung system associated with Eq.~\eqref{eq:schrodinger_bridge_appendix} \cite{Schrodinger1926,Madelung1927}.
    
    \subsection{SST reading of the bridge}
        The conservative SST interpretation is:
        \begin{align}
            \rho_\psi
            &\leftrightarrow
            \text{coarse-grained mode-support density},\\
            \theta
            &\leftrightarrow
            \text{effective circulation phase},\\
            \mathbf{u}_\psi
            &\leftrightarrow
            \text{phase-transport velocity of the low-energy envelope}.
        \end{align}
        
        At this level, the branch-resolved atomic sector of Section~\ref{sec:atomic} is read not as a replacement of the orthodox wavefunction, but as a candidate microscopic underpinning for the effective fields \(\rho_\psi\) and \(\theta\).
    
    \subsection{Two-component branch bookkeeping}
        A minimal two-branch extension consistent with the \(\sigma=\pm1\) structure of the atomic sector is
        \begin{align}
            \Psi
            =
            \begin{pmatrix}
                \psi_{+}\\
                \psi_{-}
            \end{pmatrix},
            \qquad
            \psi_{\pm}
            =
            \sqrt{\rho_{\pm}}\,e^{i\theta_{\pm}},
            \label{eq:two_component_bridge}
        \end{align}
        with total coarse-grained density
        \begin{align}
            \rho_{\mathrm{tot}} = \rho_{+}+\rho_{-}.
        \end{align}
        
        \textbf{[SPECULATIVE]} This two-component structure is the minimal SST bridge to spinor-like bookkeeping. It is not yet a first-principles derivation of fermionic spin statistics.
    
    \subsection{Topological reading of the bridge}
        The most conservative topology statement retained at appendix level is that nontrivial phase maps may encode nontrivial carrier organization. In particular, the effective phase field may carry nontrivial winding or linking data even when the coarse-grained description is written in wave language. This is the appendix-level reason why the canon treats topology and wave structure as compatible rather than as competing descriptions.
    
    \section{Dark-Sector and Galactic Continuum Law}
    \label{sec:appendix_dark_galactic}
    
    \subsection{Purpose and status}
        The main canon already retains the dark-knot taxonomy and the radial swirl pressure law. This appendix collects the next phenomenological layer: the galaxy-scale force law and the simplest rotation-curve ansatz inherited from earlier SST development.
        
        \textbf{[ORTHODOX]} The underlying force law is ordinary Euler balance in steady swirl.
        
        \textbf{[DERIVED]} SST reinterprets that balance as a candidate large-scale acceleration source.
        
        \textbf{[SPECULATIVE]} The explicit galaxy-fit profile below is a phenomenological closure ansatz, not a theorem.
    
    \subsection{Effective dark acceleration law}
        Starting from the steady radial pressure law of Eq.~\eqref{eq:swirl_pressure_law},
        \begin{align}
            \frac{1}{\rhoF}\frac{dp_{\mathrm{swirl}}}{dr}
            =
            \frac{v_\theta(r)^2}{r},
        \end{align}
        define the effective SST acceleration by
        \begin{align}
            a_{\mathrm{SST}}(r)
            :=
            \frac{1}{\rhoF}\frac{dp_{\mathrm{swirl}}}{dr}.
            \label{eq:appendix_dark_accel}
        \end{align}
        Then
        \begin{align}
            a_{\mathrm{SST}}(r)=\frac{v_\theta(r)^2}{r}.
            \label{eq:appendix_dark_accel_v2r}
        \end{align}
        
        \textbf{[DERIVED]} Equations~\eqref{eq:appendix_dark_accel}--\eqref{eq:appendix_dark_accel_v2r} are immediate consequences of the retained pressure law.
    
    \subsection{Two-component galactic rotation ansatz}
        A retained phenomenological velocity profile is
        \begin{align}
            v_{\mathrm{swirl}}(r)
            =
            \frac{C_{\mathrm{core}}}{\sqrt{1+\bigl(r_c^{\mathrm{gal}}/r\bigr)^2}}
            +
            C_{\mathrm{tail}}
            \left(1-e^{-r/r_c^{\mathrm{gal}}}\right),
            \label{eq:galactic_velocity_ansatz}
        \end{align}
        with corresponding total circular speed
        \begin{align}
            v_{\mathrm{tot}}^2(r)
            =
            v_{\mathrm{bar}}^2(r)
            +
            v_{\mathrm{swirl}}^2(r).
            \label{eq:total_circular_speed_appendix}
        \end{align}
        
        \textbf{[SPECULATIVE]} Equation~\eqref{eq:galactic_velocity_ansatz} is a phenomenological fit ansatz and should be used only at the appendix / research-track level.
    
    \subsection{Relation to dark-knot taxonomy}
        The taxonomy of Eq.~\eqref{eq:dark_knot_taxonomy} classifies low-helicity, low-instability knot sectors as dark or quasi-dark candidates. The present galactic law does \emph{not} assume that all galactic anomalies are produced by localized dark knots alone. Rather, it provides the complementary continuum statement: long-range swirl pressure gradients may produce effective extra acceleration even when the underlying microscopic carriers are topologically sparse or weakly luminous.
    
    \subsection{Use policy}
        The galactic law belongs in the canon only as:
        \begin{enumerate}
            \item an appendix-level phenomenological sector,
            \item a separate astrophysics / dark-sector track,
            \item a fit target for future rotation-curve studies.
        \end{enumerate}
        It is deliberately not promoted in v0.8.1 to the same status as the primitive constant chain, the Swirl-Clock, or the delay-selection sector.
        
% ---------------------------
% RESEARCH TRACKS for future
% ---------------------------
        \newpage
        %! Author = Omar Iskandarani
%! Companion to Swirl-String-Theory Canon-v0.8.1 (included from \texttt{SST\_CANON-v0.8.1.tex} before the manual bibliography).

\section{Research-track sectors migrated from the main canon}

\noindent
This companion document collects those sectors that remain useful for continuity and future development but are too long, too provisional, or too phenomenology-specific for the core body of \textit{Swirl-String-Theory\_Canon-v0.8.1}. Their migration out of the main canon is editorial rather than negative: each subsection remains part of the broader SST research program, but none is required for the minimal logical closure of the core canon.

\medskip
\noindent
\textbf{Editorial note (v0.8.1):} the maximal swirl tension sector ($F_{\rm swirl}^{\max}$, Rydberg bridge identities, Gibbons-class maximum-tension comparison, and the threefold interpretation table) is developed in the core document under \verb|SST_CANON-v0.8.1.tex|, subsection \texttt{Maximal Swirl Tension} (label \texttt{sec:Fmax}); it is not duplicated here.

\medskip

\subsection{Thermodynamic Radiation Interface}
A recurring legacy idea is that a local swirl-energy density may define an effective temperature field and thereby an emergent radiation sector. Let
\begin{align}
    T_{\mathrm{eff}} = T_{\mathrm{eff}}\big(\rhoE,\, \omega,\, \SwirlClock\big)
\end{align}
denote an unspecified constitutive relation. A minimal legacy interface then asks whether Planck-type spectra can emerge from a medium whose local excitation scale is set by $T_{\mathrm{eff}}$.

\textbf{[ORTHODOX]}: blackbody spectra are governed by Planck's law once a valid temperature variable is defined.

\textbf{[SPECULATIVE]}: SST does not yet derive the constitutive relation above, so this sector remains a research program rather than canonically closed physics.

\subsection{Minimal Coupling Interface to QED}
Early research notes repeatedly explored whether an effective vector potential could be built from circulation variables. The most conservative placeholder takes the form
\begin{align}
    \nabla \rightarrow \nabla - i\frac{m}{\hbar}A_{\mathrm{eff}},
    \qquad
    A_{\mathrm{eff}} = \chi \, v,
\end{align}
where $\chi$ is a scalar weight and $v$ is the local carrier velocity field.

\textbf{[ORTHODOX]}: minimal coupling is the standard gateway from free to gauge-coupled wave dynamics.

\textbf{[SPECULATIVE]}: the above substitution is only an analogy template unless a full gauge-covariant derivation is supplied.

\subsection{Ribbon-Invariant Chirality Refinement}
The older canon layers also suggested using the Calugareanu--White--Fuller relation
\begin{align}
    Lk = Tw + Wr
\end{align}
as a refinement of the knot-based matter/antimatter dictionary.

\textbf{[ORTHODOX]}: this is a standard geometric identity for framed curves and ribbons.

\textbf{[DERIVED]}: it can be used internally to separate writhe-dominated from twist-dominated realizations of the same coarse knot class.

\textbf{[SPECULATIVE]}: the mapping from $\operatorname{sign}(Tw)$ or related chirality measures to particle/antiparticle sectors remains a model hypothesis.

\subsection{Particle-Candidate Mapping Layer}
\textbf{[SPECULATIVE] Working dictionary:} a minimal knot/composition mapping retained from earlier SST work is summarized in Table~\ref{tab:particle-candidates}. This table is organizational rather than definitive; it records the current topological assignments used by the mass and taxonomy program.

\begin{table}[h]
\centering
\begin{tabular}{lll}
\hline
Knot class / composition & Candidate & Status / role \\
\hline
$3_1$ & electron & baseline lepton calibration \\
$T(2,5)$ or equivalent higher torus class & muon & lepton hierarchy candidate \\
$T(2,7)$ or equivalent higher torus class & tau & lepton hierarchy candidate \\
$5_2 + 5_2 + 6_1$ & proton & composite baryon candidate \\
analogous nearby composite class & neutron & composite baryon candidate \\
\hline
\end{tabular}
\caption{Minimal knot-class to particle-candidate summary retained as a research-track dictionary.}
\label{tab:particle-candidates}
\end{table}

\textbf{[CRITICAL NOTE]}: Table~\ref{tab:particle-candidates} is a working dictionary only. A full canonical assignment requires an explicit invariant set and a mass-functional comparison program.

\subsection{Hydrodynamic Exchange and the Pauli Barrier}
\textbf{[ORTHODOX]}: for two thin filamentary loops $C_1$ and $C_2$ in an incompressible medium, the Biot--Savart interaction energy is
\begin{align}
    E_{\mathrm{int}}
    =
    \frac{\rhoF\Gamma_1\Gamma_2}{4\pi}
    \oint_{C_1}\oint_{C_2}
    \frac{d\ell_1\cdot d\ell_2}{\norm{r_1-r_2}}.
\end{align}

\textbf{[DERIVED]}: if two identical electron-candidate loops are forced into the same spatial state, regularization at the core radius $\rc$ yields the overlap penalty
\begin{align}
    V_{\mathrm{Pauli}}(d)
    \approx
    \frac{\rhoF\Gamma_0^2}{4\pi}
    \iint
    \frac{d\ell_1\cdot d\ell_2}{\sqrt{\norm{r_1(\theta_1)-r_2(\theta_2)+d}^2+\rc^2}},
\end{align}
with maximal-overlap scale
\begin{align}
    V_{\mathrm{Pauli}}^{\max}
    \approx
    \frac{\rhoF\Gamma_0^2}{4\pi}\left(\frac{L}{\rc}\right)\mathcal{O}(1).
\end{align}

Using the canonical values
\begin{align}
    \rhoF &= 7.0\times 10^{-7}\ \mathrm{kg\,m^{-3}}, \\
    \rc &= 1.40897017\times 10^{-15}\ \mathrm{m}, \\
    \Gamma_0 &= 2\pi \rc\,\vnorm \approx 9.6836\times 10^{-9}\ \mathrm{m^2\,s^{-1}},
\end{align}
and $L\sim 2\pi a_0$ at the hydrogenic scale, one obtains the benchmark estimate
\begin{align}
    V_{\mathrm{Pauli}}^{\max} \sim 1.23\times 10^{-18}\ \mathrm{J} \sim 7.6\ \mathrm{eV}.
\end{align}

\textbf{[CRITICAL NOTE]}: this is retained as a canonical scale benchmark, not yet as a closed theorem. Its role is to show that the exclusion-energy estimate falls in an atomically relevant window rather than at obviously pathological scales.

\subsection{Numerical Benchmark Layer}
\textbf{[DERIVED] Computational-program requirement:} earlier SST work contained explicit benchmark tables for particle and atomic masses. In v0.8.1, the benchmark layer is retained as a reproducibility requirement rather than as a stand-alone authority claim.

\begin{table}[h]
\centering
\begin{tabular}{lllll}
\hline
Particle & $m_{\mathrm{exp}}$ (MeV) & $m_{\mathrm{SST}}$ (MeV) & rel. error & mode class \\
\hline
e & 0.51099895 & 0.51099895 & 0 & predictive/closure \\
$\mu$ & 105.6583745 & \dots & \dots & predictive candidate \\
$\tau$ & 1776.86 & \dots & \dots & predictive candidate \\
p & 938.272088 & \dots & \dots & predictive or closure \\
n & 939.565420 & \dots & \dots & predictive or closure \\
\hline
\end{tabular}
\caption{Benchmark summary structure to be populated from canonical scripts and archived CSV output.}
\end{table}

\textbf{[CRITICAL NOTE]}: a future v0.8.x benchmark appendix should include the exact script or package path used for generation, archived CSV outputs, mode definitions (predictive vs. exact-closure), and uncertainty propagation on canonical constants and geometric factors.

\subsection{Helicity-by-base archive (diagnostic layer)}
\textbf{[DERIVED] Archive role:} the SST helicity-by-base outputs are retained as a measurable balance diagnostic for taxonomy support, not as an independent law.
Define the midpoint-distance diagnostic
\begin{align}
    \Delta_H(K)=\abs{\widehat{\mathcal{H}}_{b}(K)+\frac{1}{2}},
\end{align}
where $\widehat{\mathcal{H}}_{b}(K)$ is the archived normalized helicity-like base indicator for class $K$.

Near-balanced values with $\Delta_H(K)\ll 1$ indicate candidate dark or quasi-dark sectors only when symmetry and instability filters are also satisfied.

\textbf{[CRITICAL NOTE]}: no single helicity-like scalar is sufficient for sector assignment; classifier decisions remain conjunction rules over symmetry, helicity balance, and low-order stability content.

\textbf{[CANON WORKFLOW NOTE]}: if SST-labelled and parallel VAM-labelled helicity outputs are numerically identical, the SST-labelled archive is the only required source for canon workflows.

\subsection{Trefoil-closure archive (benchmark discipline)}
\textbf{[DERIVED] Benchmark role:} ideal-trefoil robustness sweeps, final-best-estimate tables, and theorem-style reports are retained as benchmark-discipline infrastructure.

\textbf{[CRITICAL NOTE]}: practical best-estimate closure and theorem-level closure are distinct.
A result may remain useful for calibration while still failing strict theorem-level tail or postcheck requirements.

\textbf{[REQUIRED METADATA]}: benchmark reports should include explicit pass/fail logic, postchecks, extrapolation summaries, and local-curve diagnostics so that closure claims are auditable.

\subsection{\texorpdfstring{$\phi_{3_1}$}{phi3_1} transform archive (deferred)}
\textbf{[SPECULATIVE] Status:} retained as research-track spectral evidence for later zeta/spectral programs.

\textbf{[SCOPE RULE]}: this archive is not part of the current dark-sector or galactic canonization closure path in v0.8.x.

\subsection{Gauge-Sector Roadmap}
\textbf{[SPECULATIVE] Structural roadmap:} a retained minimal statement of the gauge program is
\begin{align}
    g_{\mathrm{eff}} = su(3) \oplus su(2) \oplus u(1),
\end{align}
with effective gauge phases organized as holonomy data of multi-director transport.

A minimal phase-based representation is
\begin{align}
    A^a_{\mu}(x) = \partial_{\mu}\theta^a(x),
\end{align}
and clock coupling acts as a common phase reference,
\begin{align}
    \theta^a(x) \mapsto \theta^a(x)+q^a\chi(x)
    \quad\Rightarrow\quad
    A^a_{\mu}\mapsto A^a_{\mu}+q^a\partial_{\mu}\chi.
\end{align}

\textbf{[CRITICAL NOTE]}: the gauge layer is retained only as a roadmap statement. Running couplings, anomaly cancellation, and realistic mixing matrices remain open work.

\section{Dark-sector and galactic canonization execution package (v0.8.x)}
\label{sec:dark-sector-execution-package}

\subsection{Stage-0 freeze: topology toolkit and benchmark protocol}
\textbf{[DERIVED] Scope freeze:} the mandatory topology toolkit for v0.8.x canonization is fixed to:
\begin{enumerate}
    \item Gauss linking integral,
    \item helicity in ideal fluids,
    \item thin-tube helicity reduction,
    \item Hopf invariant,
    \item basic torus-knot invariants $(c,b,g)$,
    \item Calugareanu--White--Fuller relation $Lk=Tw+Wr$.
\end{enumerate}

\textbf{[DERIVED] Notation freeze:} one notation family is used in this package:
\[
\rhoF,\quad v_\theta(r),\quad p_{\mathrm{swirl}}(r),\quad
\mathcal{H}(K),\quad N_u^{(1)}(K),\quad \Gamma,\quad \kappa,\quad \SwirlClock.
\]
No parallel symbol aliases are introduced for the same quantity.

\textbf{[DERIVED] Mandatory benchmark protocol template:}
\begin{enumerate}
    \item \textbf{Inputs:} canonical constants, topology identifiers, profile parameters, and dataset/version tags.
    \item \textbf{Procedure:} exact script or package path, run command, and branch/commit hash.
    \item \textbf{Outputs:} archived CSV files with fixed column schema and an attached metadata note.
    \item \textbf{Uncertainty:} first-order propagation for calibrated constants and sensitivity sweep on profile parameters.
    \item \textbf{Rerun:} deterministic rerun steps and acceptance tolerances.
\end{enumerate}

\textbf{[CANON WORKFLOW NOTE]}: every promoted quantitative claim in this package is invalid without an explicit input-to-CSV trace.

\subsection{Stage-1 freeze: dark-sector Euler pressure anchor}
\textbf{[ORTHODOX] Euler radial balance (steady, axisymmetric, azimuthal-dominant):}
\begin{align}
    \frac{1}{\rhoF}\frac{dp_{\mathrm{swirl}}}{dr}
    =
    \frac{v_\theta(r)^2}{r}.
    \label{eq:exec_swirl_pressure}
\end{align}

\textbf{[DERIVED] Flat-tail limit (constant tail speed):}
\begin{align}
    v_\theta(r)\to v_0
    \quad\Longrightarrow\quad
    p_{\mathrm{swirl}}(r)
    =
    p_0 + \rhoF v_0^2\ln\!\left(\frac{r}{r_0}\right).
\end{align}

\textbf{[DERIVED] Assumption block:}
incompressible, inviscid, stationary, axisymmetric, azimuthal-dominant flow, with validity restricted to regions where radial drift and stratification corrections are subleading.

\textbf{[DERIVED] Dimensional check:}
\[
\left[\frac{1}{\rhoF}\frac{dp}{dr}\right]
=
\frac{\mathrm{kg\,m^{-2}\,s^{-2}}}{\mathrm{kg\,m^{-3}}}
=
\mathrm{m\,s^{-2}}
=
\left[\frac{v^2}{r}\right].
\]

\textbf{[SPECULATIVE] SST interpretation:}
Eq.~\eqref{eq:exec_swirl_pressure} is interpreted as a classifier/driver layer for dark-sector support. This interpretation is model-level and separate from the orthodox derivation.

\textbf{[DERIVED] Attached observables/falsifiers:}
\begin{enumerate}
    \item Rotation-curve consistency after pressure reconstruction from archived $v_\theta(r)$ data.
    \item Null falsifier: systematic mismatch in the reconstructed pressure gradient sign or scale across validated radial windows.
\end{enumerate}

\subsection{Stage-2 freeze: measurable dark taxonomy}
\textbf{[DERIVED] Minimal classifier variables:}
symmetry class (including amphichirality status), chirality sign/state, helicity-balance indicator $\Delta_H(K)$, and low-order instability count $N_u^{(1)}(K)$.

\textbf{[DERIVED] Decision rules (symmetry-first, then stability-filtered):}
\begin{align}
    \text{Dark}
    &:\ \text{amphichiral candidate}
    \ \wedge\
    \Delta_H(K)\ll 1
    \ \wedge\
    N_u^{(1)}(K)=0,\\
    \text{Quasi-dark}
    &:\ \text{near-balanced symmetry/helicity}
    \ \wedge\
    0 < N_u^{(1)}(K)\le 2,\\
    \text{Visible}
    &:\ \text{chiral-dominant and/or helicity-unbalanced sectors}.
\end{align}

\textbf{[DERIVED] Ambiguity tie-break:}
if classes overlap, choose the class with the strongest conjunction score in the order
\[
\text{symmetry gate} \rightarrow \text{helicity balance gate} \rightarrow \text{stability gate}.
\]

\textbf{[DERIVED] Baseline templates and falsifiers:}
\begin{itemize}
    \item \textbf{Visible baseline:} trefoil-like chiral template; falsifier = persistent amphichiral-plus-balanced signature under validated data updates.
    \item \textbf{Quasi-dark baseline:} Borromean-type weak-coupling template; falsifier = stable zero-instability assignment across independent scans.
    \item \textbf{Dark baseline:} figure-eight/amphichiral template; falsifier = robust nonzero chirality or high-instability assignment.
\end{itemize}

\textbf{[CANON WORKFLOW NOTE]}:
amphichiral $\Rightarrow$ dark remains a classifier proposal tested against observables, not a closed theorem.

\subsection{Stage-3 freeze: galactic law branch gate and closure route}
\textbf{[DERIVED] Branch-gate decision:} Branch A (legacy profile refine) is selected for current v0.8.x closure because it preserves tested limits while allowing benchmark traceability.

Working profile:
\begin{align}
    v_\theta(r)
    =
    \frac{C_{\mathrm{core}}}{\sqrt{1+(r_c/r)^2}}
    +
    C_{\mathrm{tail}}\left(1-e^{-r/r_c}\right).
\end{align}

\textbf{[DERIVED] Required checks:}
small-$r$ behavior, flat-tail behavior, parameter semantics, and pressure reconstruction via Eq.~\eqref{eq:exec_swirl_pressure}.

\textbf{[DERIVED] Failure modes:}
\begin{enumerate}
    \item non-physical parameter combinations (wrong sign or unstable growth),
    \item failure to satisfy tail-limit behavior in benchmark windows,
    \item inability to reconstruct a consistent pressure profile from the same archived inputs.
\end{enumerate}

\subsection{Stage-4 freeze: probe falsification map}
\textbf{[DERIVED]} probe classes are fixed as astrophysical, laboratory pressure/swirl analog, HV-coil/thrust-inspired, optical/torsional coupling, and null dark-sector tests.

\begin{table}[h]
\centering
\begin{tabular}{p{2.7cm}p{3.1cm}p{2.5cm}p{2.2cm}p{2.2cm}}
\hline
Probe class & Equation under test & Null-result meaning & Primary confounder & Expected sign/trend \\
\hline
Astrophysical rotation curves & Eq.~\eqref{eq:exec_swirl_pressure} + profile closure & weakens pressure-law applicability window & baryonic-model bias & reconstructed $dp/dr$ consistent with $v_\theta^2/r$ \\
Lab swirl analog & Euler radial balance scaling & weakens scale-transfer claim & viscosity/edge forcing & monotone pressure rise with imposed swirl \\
HV-coil/thrust-inspired & circulation-pressure bridge claims & kills quantized-bridge claim if no scaling & EM/thermal artifacts & response tracks circulation proxy \\
Optical/torsional coupling & swirl-clock bridge signatures & weakens coupling interpretation & detector drift/noise & phase/timing shift with controlled swirl proxy \\
Null dark tests & taxonomy conjunction rules & kills classifier branch if conjunction fails & mislabelled topology class & class score decreases under contradictory diagnostics \\
\hline
\end{tabular}
\caption{Frozen equation-to-probe map for falsification-first workflow in the v0.8.x dark-sector program.}
\label{tab:exec-probe-map}
\end{table}

\textbf{[CANON WORKFLOW NOTE]}: probe outcomes can support or reject already-stated equations, but they do not replace derivation-level closure.

\subsection{Promotion snapshot after execution}
\textbf{[DERIVED]} post-freeze status in this package:
\begin{itemize}
    \item Item 6 (topology toolkit): CANON-ready in scope and notation.
    \item Item 5 (benchmark protocol): CANON-CANDIDATE pending repository-wide protocol reuse.
    \item Item 1 (Euler pressure anchor): CANON-CANDIDATE with explicit assumptions, limits, and falsifier path.
    \item Item 2 (dark taxonomy): BRIDGE with measurable classifier gates and falsifiers.
    \item Item 3 (galactic swirl law): BRIDGE with branch gate closed and benchmark route fixed.
    \item Item 4 (probe map): CANON-BRIDGE with equation/null-result mapping frozen.
\end{itemize}

        
% ---------------------------
% Bibliography
% ---------------------------
        \begin{thebibliography}{99}
            
            \bibitem{SST_delay_basic}
            J.~K. Hale and S.~M. Verduyn Lunel,
            \newblock \emph{Introduction to Functional Differential Equations},
            \newblock Springer (1993).
            
            \bibitem{CMS2026Wmass}
            CMS Collaboration,
            \newblock High-precision measurement of the W boson mass with the CMS experiment,
            \newblock Nature \textbf{652} (2026), 321--337,
            \newblock DOI: 10.1038/s41586-026-10168-5.
            
            \bibitem{PDG2024}
            S.~Navas et al. (Particle Data Group),
            \newblock Review of Particle Physics,
            \newblock Phys. Rev. D \textbf{110} (2024), 030001,
            \newblock DOI: 10.1103/PhysRevD.110.030001.
            
            \bibitem{deBlas2022}
            J.~de Blas et al.,
            \newblock Global analysis of electroweak data in the Standard Model,
            \newblock Phys. Rev. D \textbf{106} (2022), 033003,
            \newblock DOI: 10.1103/PhysRevD.106.033003.
            
            \bibitem{CDF2022Wmass}
            CDF Collaboration,
            \newblock High-precision measurement of the W boson mass with the CDF II detector,
            \newblock Science \textbf{376} (2022), 170--176,
            \newblock DOI: 10.1126/science.abk1781.
            
            \bibitem{Batchelor1967}
            G.~K. Batchelor,
            \newblock \emph{An Introduction to Fluid Dynamics},
            \newblock Cambridge University Press (1967).

            \bibitem{Jacobson2008}
            T.~Jacobson,
            \newblock Einstein-\AE ther gravity: a status report,
            \newblock PoS QG-PH (2007) 020,
            \newblock DOI: 10.22323/1.043.0020.

            \bibitem{BlasPujolasSibiryakov2011}
            D.~Blas, O.~Pujol\`as, and S.~Sibiryakov,
            \newblock Models of non-relativistic quantum gravity: the good, the bad and the healthy,
            \newblock JHEP \textbf{04} (2011) 018,
            \newblock DOI: 10.1007/JHEP04(2011)018.

            \bibitem{Monitor2017}
            B.~P. Abbott et al. (LIGO Scientific Collaboration and Virgo Collaboration),
            \newblock GW170817: Observation of gravitational waves from a binary neutron star inspiral,
            \newblock Phys. Rev. Lett. \textbf{119} (2017) 161101,
            \newblock DOI: 10.1103/PhysRevLett.119.161101.

            \bibitem{Unruh1981}
            W.~G. Unruh,
            \newblock Experimental black-hole evaporation?,
            \newblock Phys. Rev. Lett. \textbf{46} (1981), 1351--1353,
            \newblock DOI: 10.1103/PhysRevLett.46.1351.

            \bibitem{Visser1998}
            M. Visser,
            \newblock Acoustic black holes: horizons, ergospheres and Hawking radiation,
            \newblock Class. Quantum Grav. \textbf{15} (1998), 1767--1791,
            \newblock DOI: 10.1088/0264-9381/15/6/024.

            \bibitem{Painleve1921}
            P. Painlev\'e,
            \newblock La m\'ecanique classique et la th\'eorie de la relativit\'e,
            \newblock C. R. Acad. Sci. (Paris) \textbf{173} (1921), 677--680.

            \bibitem{Gullstrand1922}
            A. Gullstrand,
            \newblock Allgemeine L\"osung des statischen Eink\"orperproblems in der Einsteinschen Gravitationstheorie,
            \newblock Arkiv. Mat. Astron. Fys. \textbf{16}(8) (1922), 1--15.

            \bibitem{Schrodinger1926}
            E. Schr\"odinger,
            \newblock Quantisierung als Eigenwertproblem,
            \newblock Annalen der Physik \textbf{384} (1926), 361--376,
            \newblock DOI: 10.1002/andp.19263840404.

            \bibitem{Madelung1927}
            E. Madelung,
            \newblock Quantentheorie in hydrodynamischer Form,
            \newblock Z. Phys. \textbf{40} (1927), 322--326,
            \newblock DOI: 10.1007/BF01400372.

            \bibitem{Jackson1999}
            J.~D. Jackson,
            \newblock \emph{Classical Electrodynamics}, 3rd ed.,
            \newblock Wiley (1999).

            \bibitem{Planck1901}
            M. Planck,
            \newblock Ueber das Gesetz der Energieverteilung im Normalspectrum,
            \newblock Annalen der Physik \textbf{309} (1901), 553--563,
            \newblock DOI: 10.1002/andp.19013090310.

            \bibitem{Wien1896}
            W. Wien,
            \newblock Ueber die Energieverteilung im Emissionsspectrum eines schwarzen K\"orpers,
            \newblock Annalen der Physik \textbf{294} (1896), 662--669,
            \newblock DOI: 10.1002/andp.18962940310.

            \bibitem{Calugareanu1961}
            G. C\u{a}lug\u{a}reanu,
            \newblock Sur les classes d'isotopie des noeuds tridimensionnels et leurs invariants,
            \newblock Czech. Math. J. \textbf{11} (1961), 588--625.

            \bibitem{White1969}
            J.~H. White,
            \newblock Self-linking and the Gauss integral in higher dimensions,
            \newblock Amer. J. Math. \textbf{91} (1969), 693--728,
            \newblock DOI: 10.2307/2373348.

            \bibitem{Saffman1992}
            P.~G. Saffman,
            \newblock \emph{Vortex Dynamics},
            \newblock Cambridge University Press (1992).

            \bibitem{Moffatt1969}
            H.~K. Moffatt,
            \newblock The degree of knottedness of tangled vortex lines,
            \newblock J. Fluid Mech. \textbf{35} (1969), 117--129,
            \newblock DOI: 10.1017/S0022112069000991.

            \bibitem{Zurek2003}
            W.~H. Zurek,
            \newblock Decoherence, einselection, and the quantum origins of the classical,
            \newblock Rev. Mod. Phys. \textbf{75} (2003), 715--775,
            \newblock DOI: 10.1103/RevModPhys.75.715.

            \bibitem{Weinberg1967}
            S.~Weinberg,
            \newblock A model of leptons,
            \newblock Phys. Rev. Lett. \textbf{19} (1967), 1264--1266,
            \newblock DOI: 10.1103/PhysRevLett.19.1264.

            \bibitem{Gibbons2002}
            G.~W. Gibbons,
            \newblock The maximum tension principle in general relativity,
            \newblock Found. Phys. \textbf{32} (2002), 1891--1901,
            \newblock DOI: 10.1023/A:1022370717626.

        \end{thebibliography}

\end{document}